%% =====================================================================================================================
%% Sustainable Local RAG Architectures for Energy Governance:
%% A Green AI Approach via Document Slice Contextualization
%%
%% Springer Nature journal manuscript, sn-jnl class (Math & Physical Sciences, numbered references).
%% Compile with the Springer Nature LaTeX template package (sn-jnl.cls, sn-mathphys-num.bst) and refs.bib.
%%
%% DRAFT STATUS. This file was drafted with AI assistance (Claude, Anthropic) from the CITDS conference
%% manuscript (citds.tex) and the repository's dynamic-router ablation study. Conventions:
%%   \todo{...}    red text: missing figure, unconfirmed measurement, provisional statistic or missing citation
%%   \verify{...}  red text: how a measured number was computed and where in the repository it can be checked
%% Both must be resolved and removed before submission (coloured text is not reproduced in the publisher's XML).
%% Figures are placeholders only: no image files are referenced.
%% The scripts that recompute the reported numbers are in RnD/verification/ (see its README.md); \verify notes name them.
%% =====================================================================================================================
\documentclass[pdflatex,sn-mathphys-num]{sn-jnl}

\usepackage{graphicx}
\usepackage{multirow}
\usepackage{amsmath,amssymb,amsfonts}
\usepackage{amsthm}
\usepackage{mathrsfs}
\usepackage[title]{appendix}
\usepackage{xcolor}
\usepackage{textcomp}
\usepackage{manyfoot}
\usepackage{booktabs}
\usepackage{algorithm}
\usepackage{algorithmicx}
\usepackage{algpseudocode}
\usepackage{listings}
\usepackage{microtype}

%% --- placeholder / provenance macros (remove before submission) -----------------------------------------------------
\newcommand{\todo}[1]{\textcolor{red}{#1}}
\newcommand{\verify}[1]{\todo{[VERIFY: #1]}}

%% --- notation -------------------------------------------------------------------------------------------------------
\newcommand{\gco}{g\,CO$_2$eq}
\newcommand{\Rec}[1]{Recall@#1}
\newcommand{\MRR}[1]{MRR@#1}
\newcommand{\MAR}[1]{MAR@#1}
\newcommand{\Slice}{\operatorname{Slice}}
\newcommand{\rnk}{\operatorname{rank}}

\lstset{basicstyle=\footnotesize\ttfamily, breaklines=true, breakatwhitespace=true, frame=single, columns=fullflexible, keepspaces=true, xleftmargin=2pt, framexleftmargin=2pt}

\theoremstyle{thmstyleone}%
\newtheorem{theorem}{Theorem}%
\newtheorem{proposition}[theorem]{Proposition}%
\theoremstyle{thmstylethree}%
\newtheorem{definition}{Definition}%

\raggedbottom

\begin{document}

\title[Sustainable Local RAG for Energy Governance]{Sustainable Local RAG Architectures for Energy Governance: A Green AI Approach via Document Slice Contextualization}

\author*[1]{\fnm{Martin Tibor} \sur{S\'andor}}\email{sandor.martin@inf.unideb.hu}
\author[2]{\fnm{R\'obert} \sur{Lakatos}}\email{lakatos.robert@inf.unideb.hu}

\affil*[1]{\orgdiv{Faculty of Informatics}, \orgname{University of Debrecen}, \orgaddress{\city{Debrecen}, \country{Hungary}}}
\affil[2]{\orgdiv{Department of Data Science and Visualization, Faculty of Informatics}, \orgname{University of Debrecen}, \orgaddress{\city{Debrecen}, \country{Hungary}}}

%% VERIFY (abstract; kept free of red text): every number below is stated again with its provenance note in the body:
%%   1.37 mean radius, 32 % skipped, 0.952 vs 0.954, 106.7 -> 13.2 Wh, 21.8 -> 2.8 g, 54.9 -> 8.8 min  -> Table 7 (tab:primary) and Sect. 4.2
%%   0.6 points and 71.8 % on the policy corpus                                                     -> Table 8 (tab:policy) and Sect. 4.5
\abstract{\textbf{Background:} Retrieval-augmented generation pipelines that contextualise every chunk against its whole document follow a ``Red AI'' pattern: quadratic self-attention and a linearly growing key--value cache push index-time inference beyond the 8\,GB of a consumer GPU, while energy-sector auditing and policy work, bound by data-sovereignty rules, must run locally.

\textbf{Results:} We formalise Document Slice contextualisation, which conditions each chunk summary on a bounded window of neighbouring chunks, and add a two-tier routing agent that picks the radius per chunk with a regex gate and one single-token forward pass of the same 4B-parameter model. On 25 scientific papers (382 chunks, 250 multi-hop questions) the router chose a mean radius of 1.37, skipped 32\,\% of the summaries and reached \Rec{20} of 0.952 against 0.954 for full-document contextualisation, with higher MRR, while cutting tracked generation energy from 106.7 to 13.2\,Wh ($-87.6$\,\%), emissions from 21.8 to 2.8\,\gco{} and generation time from 54.9 to 8.8\,min. On four European Commission policy documents the fixed-radius variant stayed within 0.6 points of the baseline at 71.8\,\% lower energy.

\textbf{Conclusion:} Bounded, routed contextualisation offers a PEP-AI pathway (Precise, Explainable, Provable) to decentralised compliance auditing: reproducible metrics, inspectable per-chunk decisions and per-run energy ledgers.}

\keywords{Retrieval-Augmented Generation, Green AI, Small Language Models, Edge Inference, Energy Governance, Contextual Retrieval, Carbon Footprint}

\maketitle

\section{Introduction}\label{sec:intro}

Large language models (LLMs) have changed how organisations search, summarise and reason over documents \cite{bengioetal2003,zhu2025informationretrieval}, and they have concentrated capability among those who can afford data-centre-scale compute: the ``compute divide'' separates institutions that can train and serve frontier models from those that cannot \cite{besiroglu2024computedividemachinelearning,sevilla2022compute,thompson2020computational,hooker2021hardware}. Schwartz et al.\ named the underlying culture ``Red AI'', in which accuracy is bought with exponentially growing compute, and contrasted it with ``Green AI'', which reports efficiency as a first-class result \cite{schwartz2020green}. The energy, carbon and water footprints of training and serving such models are documented in detail \cite{strubell2019energycost,patterson2021carbon,Luccioni_2024,li2025makingaithirstyuncovering,devries2023growing,wu2022sustainable}, and reporting frameworks exist to make them comparable \cite{henderson2020systematic,lacoste2019quantifyingcarbonemissionsmachine,lannelongue2021green,tabbakh2024GreenAI}.

Energy governance is a domain in which these costs cannot simply be externalised to a cloud provider. Regulatory auditing of grid and market operators, the transposition of the recast Energy Efficiency Directive \cite{eueed2023}, the reporting duties of the Energy Union governance framework \cite{eugovernance2018} and the Commission's plan for digitalising the energy system \cite{eu2022digitalising} all produce long, structured documents that decision makers must query precisely and defensibly. At the same time the General Data Protection Regulation \cite{gdpr2016} and the Artificial Intelligence Act \cite{euaiact2024} impose data-sovereignty, transparency and record-keeping obligations on automated decision support \todo{[CITATION: AI-assisted decision support in energy regulation or policy analysis]}. For a public authority or a small operator, on-premises retrieval on commodity hardware is therefore a compliance property rather than a preference, and it has to fit the 8\,GB of video memory of a typical workstation GPU \cite{zhou2019edge}.

Retrieval-augmented generation (RAG) grounds model answers in retrieved text \cite{lewisetal2021,info16090766,ram2023incontext}, and production systems now combine dense retrieval \cite{karpukhin2020dpr} with lexical scoring such as BM25 \cite{robertson1995okapi,robertson2009bm25} and late-interaction re-ranking \cite{khattab2020colbert,santhanam2022colbertv2,kiran2025hybrid,wang2024searching}. Their weakest link is chunking: a passage cut from its document loses the referents that make it retrievable \cite{merola2025reconstructingcontextevaluatingadvanced,smith2024evaluating,qu2024semanticchunkingworthcomputational,chen2023densex}, and layout-aware segmentation only partly repairs this \cite{auer2024doclingtechnicalreport,verma2025s2chunkinghybridframework}. Context-aware representations have been proposed at embedding time \cite{gunther2024latechunking} and, most influentially, at index time, by prepending a model-written summary of the chunk's role before embedding \cite{anthropic2024contextual}.

Index-time contextualisation as practised industrially sends the whole document with every chunk. For a document of $n$ chunks and $L$ tokens this costs $(n+1)L$ prompt tokens per document, and because self-attention scales quadratically in sequence length \cite{vaswani2017attention,tay2022efficient,kaplan2020scaling} the attention work grows with $\sum_i (L+\ell_i)^2$. The memory consequence is more decisive than the arithmetic: the key--value cache of a decoder grows linearly with the context window \cite{pope2023efficiently,kwon2023vllm,ainslie2023gqa}, and once weights plus cache exceed physical VRAM the runtime pages to host memory over PCIe, where throughput collapses \cite{sheng2023flexgen}. Efficient kernels \cite{dao2022flashattention}, sparse attention \cite{beltagy2020longformer} and quantisation \cite{jacob2018quantization} lower the constants, not the dependence on $L$. On an 8\,GB GPU the full-document design thus reproduces the Red AI regime in miniature \cite{shen2024RAGDesignTradeoffs,zhao2025RAGDesignDecisions,samsi2023words}.

Long contexts are not only expensive; they are used badly. Liu et al.\ showed that models retrieve information best from the beginning and the end of a prompt and degrade for the middle \cite{liu2023lostmiddlelanguagemodels}; the effect is architectural rather than model-specific \cite{guo2024serialpositioneffectslarge,menschikov2025earlytokenbiasmodelspecificlanguagespecific,hsieh2024found}, it persists in embedding models \cite{lee2025quantifyingpositionalbiasestext} and in multimodal retrieval \cite{yao2025spotlighthiddenbiasundermining}. A summariser handed a whole paper to describe one paragraph thus pays quadratically for context it largely ignores.

Context-expansion strategies such as sentence-window and parent--child retrieval \cite{gao2024retrievalaugmentedgenerationlargelanguage,Liu_LlamaIndex_2022}, recursive summary trees \cite{sarthi2024raptorrecursiveabstractiveprocessing} and graph-based summarisation \cite{edge2024graphrag} enrich chunks with neighbouring or hierarchical context, but with variable-length inputs that are hard to bound. Adaptive systems decide per query how much retrieval to perform \cite{jeong2024adaptiverag,asai2024selfrag} or which model to call \cite{ong2024routellm,chen2023frugalgpt}. Missing is a per-chunk decision, taken at index time, of how much neighbouring context a chunk needs, under a hard bound that keeps every prompt inside a fixed memory envelope.

This article extends our conference study of Document Slice contextualisation \cite{sandor2025documentslicecontextualretrieval} in that direction. Its contributions are: (i) a formal definition of the bounded, boundary-clamped slice and of the index-time cost it implies; (ii) a two-tier routing agent that assigns each chunk a radius $k\in\{0,2\}$ with a deterministic table gate and a single-token forward pass of the same 4-billion-parameter model that writes the summaries, choosing a mean radius of 1.37 and skipping a third of the chunks; (iii) a complexity and memory model that explains why the full-document baseline leaves the 8\,GB envelope; (iv) an evaluation on 250 multi-hop questions in which the routed pipeline reaches \Rec{20} of 0.952 against 0.954 for the full-document baseline, inside the run-to-run band of the summariser, at 12.4\,\% of the baseline's tracked generation energy; (v) a cross-domain check on European Commission policy documents; and (vi) a measurement protocol, including a noise-floor analysis and per-run energy ledgers, that we consider a prerequisite for Green AI claims. Section~\ref{sec:method} presents the method and its cost model, Sect.~\ref{sec:results} the results, Sect.~\ref{sec:discussion} the operational implications and Sect.~\ref{sec:conclusion} concludes.

\section{Methodology}\label{sec:method}

The pipeline (Fig.~\ref{fig:architecture}) runs entirely on one workstation with a single 8\,GB GPU. A PDF is parsed into layout-aware chunks (Sect.~\ref{sec:parsing}); for every chunk a routing agent decides how much neighbouring context, if any, a small language model (SLM) should see when it writes the chunk's context summary (Sects.~\ref{sec:slice} and \ref{sec:router}); the summary is prepended to the chunk, embedded and indexed for hybrid retrieval (Sect.~\ref{sec:retrieval}). Sections~\ref{sec:complexity}--\ref{sec:energy} give the cost model, the evaluation metrics, the datasets and the energy accounting. Throughout, the only GPU-resident model during context generation is \texttt{gemma3:4b-it-qat} \cite{gemma3report2025}, a quantisation-aware-trained 4-billion-parameter model served by Ollama \cite{ollama}; it acts as router and as summariser, so no second model competes for memory.

\begin{figure}[htbp]
  \centering
  \todo{[FIGURE INSERTION: block diagram of the end-to-end on-premise pipeline. Left to right: PDF $\rightarrow$ Docling layout analysis and HybridChunker (2\,000-token cap, nomic tokenizer) $\rightarrow$ chunk sequence $c_1..c_n$ $\rightarrow$ routing agent (Tier~1 regex table gate; Tier~2 single-token forward pass of gemma3:4b-it-qat) $\rightarrow$ slice builder (radius $k_i\in\{0,2\}$, boundary-clamped) $\rightarrow$ summariser (same resident gemma3:4b-it-qat runner, num\_ctx 16\,384) $\rightarrow$ ``summary + chunk'' $\rightarrow$ nomic-embed-text-v1.5 (768-d) $\rightarrow$ LanceDB (vector column + BM25 full-text index) $\rightarrow$ hybrid search $\rightarrow$ ColBERTv2 re-ranker $\rightarrow$ top-20 chunks. Draw an ``8\,GB VRAM envelope'' box around the model stages and annotate that Ollama and the PyTorch embedder never share the GPU (sequential execution). Mark the k=0 bypass (chunk embedded without a summariser call).]}
  \caption{End-to-end on-premise architecture. All stages run on one 8\,GB consumer GPU; the routing agent and the summariser share one resident \texttt{gemma3:4b-it-qat} runner with a 16\,384-token context window, and the embedding model runs after the runner has been stopped. Chunks routed to $k=0$ bypass the summariser entirely and are embedded as they are. The peak GPU memory observed over the whole pipeline was 7\,695\,MiB.}
  \label{fig:architecture}
\end{figure}

\subsection{Layout-aware document parsing}\label{sec:parsing}

\begin{definition}[Chunk sequence]\label{def:chunks}
Let $\tau$ be the tokenizer of the embedding model and $\ell_{\max}$ a token budget. A document is represented by the ordered sequence of chunks it is segmented into,
\begin{equation}\label{eq:chunks}
D=(c_1,\dots,c_n),\qquad \ell_j=|\tau(c_j)|\le \ell_{\max}=2000,\qquad L=\sum_{j=1}^{n}\ell_j,\qquad \mathrm{id}(c_j)=\mathrm{SHA256}(\tilde c_j),
\end{equation}
where $\tilde c_j$ is the cleaned text of the chunk and $L$ is the token length of the whole document. The identifier is a content hash, so the same chunk receives the same identifier in every experimental condition and the question sets can reference chunks without depending on any particular index.
\end{definition}

Segmentation is performed with Docling \cite{auer2024doclingtechnicalreport,livathinos2025doclinglayoutanalysis}: its default converter runs layout analysis, table-structure recognition and OCR (RapidOCR) on each page, and its \texttt{HybridChunker} \cite{docling2024chunkingdocumentation} first splits the document along structural elements (headings, paragraphs, list items, tables, captions) and then merges adjacent elements of the same section while the token count, measured with the tokenizer of \texttt{nomic-embed-text-v1.5}, stays below $\ell_{\max}=2000$ (\texttt{merge\_peers=True}). Unlike fixed-size token windows, the resulting chunks respect the document's hierarchy: a table, its caption and the paragraph that discusses it tend to stay together, and a chunk rarely stops mid-sentence. Each chunk text is then normalised (Unicode NFKD, typographic dashes and quotes mapped to ASCII, URLs removed, whitespace collapsed) before hashing (Fig.~\ref{fig:docling}).

Two properties of the output matter for what follows. First, Docling serialises a recognised table as one line of ``\texttt{\textless row\textgreater, \textless column\textgreater = \textless value\textgreater.}'' cells and emits captions on their own line, so tables are detectable with a regular expression and need no model call to be recognised. Second, layout artefacts survive in the text: formula placeholders (\texttt{\textless!-- formula-not-decoded --\textgreater}) in 51 of the 382 chunks and undecodable font glyph codes (\texttt{GLYPH\textless..\textgreater}) in 93 of them. Table~\ref{tab:parsing} summarises the corpus that results.

\begin{figure}[htbp]
  \centering
  \todo{[FIGURE INSERTION: flow diagram of the parsing stage. Boxes: PDF page $\rightarrow$ Docling layout analysis (text blocks, tables, captions, formulas) $\rightarrow$ structural elements $\rightarrow$ HybridChunker (token-bounded merge of peers within a section, nomic tokenizer, 2\,000 tokens) $\rightarrow$ cleaning (NFKD, dash/quote normalisation, URL removal, whitespace) $\rightarrow$ SHA-256 identifier $\rightarrow$ frozen chunk JSON \{text, document, id\}. Annotate a serialised table line (``\texttt{row, column = value.}'') and a formula placeholder as examples of the artefacts the router must cope with.]}
  \caption{Layout-aware parsing and chunking. Docling's layout models produce structural elements; the hybrid chunker merges peers within a section under the 2\,000-token budget of the embedding model's tokenizer; the cleaned text is hashed into a stable identifier. On the 25-paper corpus this yields 382 chunks with a median length of 2\,090 characters, of which 42 carry a serialised table.}
  \label{fig:docling}
\end{figure}

\begin{table}[htbp]
\caption{Structural parsing metadata and chunk statistics of the scientific corpus (Docling HybridChunker output).}\label{tab:parsing}
\begin{tabular}{@{}p{5.4cm}p{4.3cm}p{3.7cm}@{}}
\toprule
Quantity & Value & Note \\
\midrule
Documents & 25 & IEEE and Nature-portfolio papers \\
Chunks & 382 & 380 distinct identifiers\footnotemark[1] \\
Chunks per document (min/median/max) & 8/14/33 & \\
Chunk length, characters (min/Q1/median/mean/Q3/max) & 4/878/2\,090/2\,766/4\,111/10\,260 & total 1\,056\,773 \\
Chunk length, tokens (min/Q1/median/mean/Q3/max) & 1/180/450/666/1\,021/1\,999 & nomic tokenizer; total 253\,488; s.d.\ 606 \\
Chunks per token band ($\le 250$/251--500/501--1\,000/1\,001--1\,500/1\,501--2\,000) & 119/88/78/37/60 & none above the 2\,000-token cap; 25 chunks at $\ge 1\,900$ \\
Characters per token & 4.17 & 1\,056\,773/253\,488 \\
Document length, characters (min/median/max) & 17\,760/43\,324/76\,161 & $\approx$ 4.3k/10.4k/18.3k tokens \\
Chunks with formula placeholders & 51 & \texttt{\textless!-- formula-not-decoded --\textgreater} \\
Chunks with glyph codes & 93 & \texttt{GLYPH\textless..\textgreater} \\
Chunks carrying a serialised table (Tier-1 gate) & 42 & cells or caption pattern \\
Token budget/merge policy & 2\,000/merge peers & \texttt{HybridChunker} defaults otherwise \\
\botrule
\end{tabular}
\footnotetext[1]{The 43-character chunk ``The authors declare no competing interests.'' occurs in three papers and therefore has one identifier; none of the three is a gold chunk.}
\end{table}

\subsection{Bounded document slices}\label{sec:slice}

\begin{definition}[Document slice]\label{def:slice}
For a chunk $c_i\in D$ and a context radius $k\in\mathbb{N}_0$ the document slice is the set of chunks within $k$ positions of $c_i$,
\begin{equation}\label{eq:slice}
\Slice(c_i,k)=\{\,c_j\in D \mid \max(1,\,i-k)\le j\le \min(n,\,i+k)\,\}.
\end{equation}
\end{definition}

Definition~\ref{def:slice} refines the chunk sequence of Definition~\ref{def:chunks} (Eq.~\eqref{eq:chunks}). Near a document boundary the window of Eq.~\eqref{eq:slice} is truncated on one side. The implementation compensates by extending the other side, so that every prompt carries the same number of chunks whenever the document is long enough (Fig.~\ref{fig:clamping}). With $u_i=\max(0,\,k-i+1)$ chunks missing before $c_i$ and $o_i=\max(0,\,i+k-n)$ chunks missing after it, the window actually sent to the model is
\begin{equation}\label{eq:window}
W_i(k)=\bigl\{\,j : \max(1,\,i-k-o_i)\le j\le \min(n,\,i+k+u_i)\,\bigr\},\qquad |W_i(k)|=\min(2k+1,\,n),
\end{equation}
and the target chunk sits in the centre of the window except at the edges, which keeps it away from the attention-degraded middle of long prompts \cite{liu2023lostmiddlelanguagemodels}. The summariser receives three messages: a fixed system prompt $\pi_s$ (Listing~\ref{lst:summariser}), a user message containing the concatenated window text under the header ``\texttt{FULL DOCUMENT:}'' (the header is kept verbatim from the full-document baseline so that the only difference between the two conditions is the content), and a user message containing the chunk under ``\texttt{CHUNK:}''. Its answer $s_i$, typically one or two sentences (about 200 characters), is prepended to the chunk, and the string $s_i \oplus \text{``\textbackslash n\textbackslash n''} \oplus c_i$ is what is embedded and indexed. The prompt length of one call, in tokens, is therefore
\begin{equation}\label{eq:prompt}
P_i(k)=s+\sum_{j\in W_i(k)}\ell_j+\ell_i \qquad\text{versus}\qquad P_i^{\mathrm{full}}=s+L+\ell_i
\end{equation}
for the full-document design, where $s$ is the length of the system prompt (789 characters, 164 tokens). Under $\ell_{\max}=2000$ every slice prompt is bounded by $s+(2k+2)\ell_{\max}$, i.e.\ by 16\,190 tokens for $k=3$ and 12\,190 for $k=2$, regardless of the document, which is why a fixed context window of 16\,384 tokens suffices for the whole corpus, whereas the full-document baseline needed 30\,000 tokens and still had to page (Sect.~\ref{sec:complexity}).

Two readings of $k=0$ occur in this article and must not be confused. In the static ablation of Sect.~\ref{sec:static} ``$k=0$'' means that the summariser is still called, with the chunk itself as its only context (a chunk-only summary). In the routed pipeline a radius of $0$ assigned by the router means that no summariser call is made at all and the chunk is embedded as it is; this is the cheaper and, as Sect.~\ref{sec:results} shows, equally accurate option for chunks that already describe themselves.

\begin{figure}[htbp]
  \centering
  \todo{[FIGURE INSERTION: schematic of Eq.~\eqref{eq:window} for a document of $n=14$ chunks and $k=2$. Three rows of 14 boxes; highlight the window for $i=1$ (chunks 1--5, shifted right by $u_1=2$), $i=7$ (chunks 5--9, centred) and $i=14$ (chunks 10--14, shifted left by $o_{14}=2$); mark the target chunk in each row. The conference figure \texttt{images/underflow\_adjusted\_doc\_slice\_alt.png} (not in the repository) shows the same mechanism and can be reused.]}
  \caption{Boundary clamping of the slice window. When the window of radius $k$ would start before the first chunk or end after the last one, the truncated side is compensated on the other side, so every prompt carries $\min(2k+1,n)$ chunks and the same number of tokens up to the chunk-length bound; the target chunk stays central except at the edges.}
  \label{fig:clamping}
\end{figure}

\begin{lstlisting}[caption={System prompt of the summariser (from \texttt{Models/chunker\_full\_doc.Modelfile}; typographic quotes and dashes rendered in ASCII).},label={lst:summariser}]
You are a summarization assistant specialized in contextual chunk analysis.

Your task is to:
- Receive a CHUNK of text and the FULL DOCUMENT it comes from.
- Output ONLY a short, concise summary of the chunk's relevance, meaning, or contribution to the document as a whole.
- DO NOT include introductions, explanations, or phrases like "This chunk is..." -- just return the clean, succinct summary.

Be objective and focus only on what the chunk contributes in the context of the full document.

Examples of good output:
- "Introduces the central thesis about climate policy reform."
- "Provides supporting evidence for the argument on economic inequality."
- "Transitions from background information to proposed methodology."

Respond only with the summary -- no preamble, no conclusion.
\end{lstlisting}

\subsection{Dynamic radius routing}\label{sec:router}

The static ablation of our conference study \cite{sandor2025documentslicecontextualretrieval} showed that $k=2$ and $k=3$ are the best fixed radii but also that a fixed radius is wasteful in two directions: chunks that already describe themselves (abstracts, conclusions, reference lists, author and funding statements, layout debris) gain nothing from a summary, and chunks that carry a serialised table are summarised \emph{worse} when the window is wide, because the model then describes the surrounding prose or degenerates into one of the examples of its own system prompt. We therefore let a routing agent $\mathcal R$ choose the radius per chunk. The agent has two tiers, both cheap, and is applied to every chunk of a document before any summary is written, so that the summariser's request sequence is independent of the routing decisions.

\subsubsection{Tier 1: deterministic table gate}
A chunk is a table chunk if one of its lines contains at least three serialised cells and is dense in cells, or if a table caption starts a line:
\begin{equation}\label{eq:gate}
\mathrm{table}(c)=\Bigl[\exists\,\text{line }\lambda\subseteq c:\ \#\mathrm{cells}(\lambda)\ge 3 \ \wedge\ |\lambda|/\#\text{``\texttt{ = }''}(\lambda)\le 100\Bigr]\ \vee\ \mathrm{caption}(c),
\end{equation}
where a cell matches the pattern ``\texttt{, \textless col\textgreater = \textless val\textgreater.}'' and a caption matches ``\texttt{TABLE I}'', ``\texttt{Table 1:}'', ``\texttt{Table 1 |}'' and their variants at the start of a line. In-text mentions (``shown in Table~1'') and mathematical prose (``$V=\{1,\dots,m\}$'') do not match. Table chunks receive $k_T=2$ without a model call.

\subsubsection{Tier 2: one forward pass of the SLM}
Every other chunk is classified by the resident \texttt{gemma3:4b-it-qat} in a single forward pass that emits one token. The prompt consists of the router system prompt $\pi_r$ (Listing~\ref{lst:router}), ten in-context exemplars $\mathcal X$ shown as ``\texttt{Passage: \ldots / Answer: S}'' pairs, and the query chunk clipped to its first 1\,500 and last 500 characters. Rather than parsing the generated text, the agent reads the log-probabilities of the top-20 candidate tokens returned by the server and takes the argmax over the two answer letters, which cannot fail to parse; if a server returns no log-probabilities the emitted token is parsed, and if that fails too the paper's default $k_{\mathrm{default}}=2$ is used. The letter map is $\kappa(\mathrm S)=0$ (self-describing: no summariser call) and $\kappa(\mathrm N)=2$ (needs its neighbours). The complete policy is
\begin{equation}\label{eq:policy}
\mathcal R(c_i)=
\begin{cases}
k_{\mathrm{fix}} & \text{if a constant radius } k_{\mathrm{fix}} \text{ is set (control runs)},\\[2pt]
k_T=2 & \text{if } \mathrm{table}(c_i) \quad\text{(Tier 1)},\\[2pt]
\kappa\Bigl(\displaystyle\arg\max_{\lambda\in\{\mathrm S,\mathrm N\}}\log p_\theta\bigl(\lambda\mid \pi_r,\mathcal X,\mathrm{clip}(c_i)\bigr)\Bigr) & \text{otherwise} \quad\text{(Tier 2)},
\end{cases}
\end{equation}
and the quantities reported for a run are the mean radius and the number of summariser calls,
\begin{equation}\label{eq:kbar}
\bar k=\frac1n\sum_{i=1}^{n}k_i,\qquad N_{\mathrm{calls}}=\bigl|\{\,i : k_i>0\,\}\bigr| .
\end{equation}
Algorithm~\ref{alg:router} lists the procedure of Eq.~\eqref{eq:policy}; Table~\ref{tab:prompts} collects the prompt templates and decision criteria, and Eq.~\eqref{eq:kbar} defines the two summary statistics reported for every run.

The routing pass is inexpensive. One forward pass takes about 0.16\,s on the RTX~4060 (median 155\,ms, range 96--287\,ms, with the ten exemplars forming a cached prefix of roughly 2\,000 tokens), so the whole 382-chunk corpus is routed in about one minute, and because the router uses the same model and the same 16\,384-token context window as the summariser, no model is unloaded or reloaded between the two phases.\verify{latencies printed by \texttt{RnD/ollama\_dynamic\_slice\_length.ipynb} (mean 157\,ms, median 155\,ms, 96--287\,ms per chunk on the 119-item gold set; 382 chunks in 60\,s); the router prompt averaged 2\,062 tokens (max 2\,602) in the same notebook; ``Summariser calls that (re)loaded the model: 0'' is printed by cell~2 of \texttt{RnD/ablation\_dynamic\_doc\_slice\_generation.ipynb}} Two design choices follow from the study reported in Sect.~\ref{sec:routing-results}: the class$\to$radius mapping matters more than the wording of the prompt (a four-class genre prompt with radii 0/1/2/3 and a two-class prompt with radii 0/2 route the benchmark almost identically once tables are gated below 3), and the exemplars must not come from the evaluation corpus. The ten exemplars are therefore taken from six papers that are not part of any benchmark (109 chunks from Lempel and Ziv 1976, SRCNN, LeCun et al.\ 2015, AlphaGo Zero, the Sycamore quantum-supremacy paper and a 2023 federated-learning survey, chunked with exactly the corpus recipe); with them the two-class router agrees with the collapsed gold labels of a 119-chunk hand-labelled set on 82 of the 99 chunks that reach Tier~2, and the table gate fires on all 20 gold table chunks with no false positive.\verify{gold set \texttt{RnD/dynamic\_slice\_length/slice\_radius\_gold.json} (119 items, 99 with tier \texttt{slm}, 20 tables); letters from \texttt{RnD/dynamic\_slice\_length/router\_experiments/results/routes\_binary\_heldout.json}; gold $k=0\to$S, $k\ge1\to$N; agreement $82/99=82.8\,\%$, confusion S$\to$N 10, N$\to$S 7; exemplar sources in \texttt{default\_few\_shots()} of \texttt{RnD/utils/dynamic\_slice\_prediction.py} and \texttt{RnD/split\_documents\_router/}} A further caveat is numerical: at temperature~0 the argmax is deterministic for given logits, but the server's logits depend slightly on whether the shared prefix was served from its cache, so chunks whose two letters are within about 0.1--0.3 nats can flip between runs; the agent exposes the letter log-probabilities so that such near-ties can be audited, which is the ``explainable'' half of the pathway we argue for in Sect.~\ref{sec:discussion}.

\begin{algorithm}[htbp]
\caption{Two-tier radius routing for one document (as implemented in \texttt{utils/dynamic\_slice\_prediction.py})}\label{alg:router}
\begin{algorithmic}[1]
\Require chunk sequence $D=(c_1,\dots,c_n)$; system prompt $\pi_r$; exemplars $\mathcal X$; letter map $\kappa$; $k_T$; $k_{\mathrm{default}}$
\Ensure radii $k_1,\dots,k_n$ and, per chunk, the tier and the letter log-probabilities
\For{$i=1$ \textbf{to} $n$}
  \If{a constant radius $k_{\mathrm{fix}}$ is configured}
    \State $k_i\gets k_{\mathrm{fix}}$; tier $\gets$ \textsc{fixed}
  \ElsIf{$\mathrm{table}(c_i)$ by Eq.~\eqref{eq:gate}}
    \State $k_i\gets k_T$; tier $\gets$ \textsc{table}
  \Else
    \State $m\gets[\pi_r;\ \{\text{``Passage: } x\ /\ \text{Answer: } \kappa^{-1}(k_x)\text{''}\}_{(x,k_x)\in\mathcal X};\ \text{``Passage: } \mathrm{clip}(c_i,1500,500)\ /\ \text{Answer:''}]$
    \State $\mathcal L\gets$ top-20 next-token log-probabilities of the SLM for $m$ \Comment{\texttt{num\_predict}=1, temperature 0}
    \State $\sigma(\lambda)\gets \mathcal L[\lambda]$ for $\lambda\in\{\mathrm S,\mathrm N\}$ present in $\mathcal L$
    \If{$\sigma=\emptyset$} \State parse the emitted token; on failure $k_i\gets k_{\mathrm{default}}$
    \Else \State $k_i\gets\kappa(\arg\max_\lambda \sigma(\lambda))$
    \EndIf
    \State tier $\gets$ \textsc{slm}
  \EndIf
\EndFor
\State \Return $(k_i,\ \text{tier}_i,\ \sigma_i)_{i=1}^{n}$ \Comment{summaries are generated afterwards, in document order, with radius $k_i$; $k_i=0$ means no call}
\end{algorithmic}
\end{algorithm}

\begin{lstlisting}[caption={System prompt of the two-class routing agent (\texttt{BINARY\_SYSTEM\_PROMPT} in \texttt{utils/dynamic\_slice\_prediction.py}).},label={lst:router}]
You are shown a passage taken from a scientific paper. Decide whether it describes itself. Answer with a single letter.

S - Self-describing. The passage already tells the reader what the paper is about or needs no such information: the abstract; a conclusion or summary of the whole paper (typically opening with 'In this paper', 'In conclusion', 'In summary', 'This paper has proposed', 'We have proposed'); or non-content matter such as a reference list, acknowledgements, funding, author biographies or contributions, competing interests, publication metadata, licence text, layout garbage. A conclusion is S even when it is full of numbers and technical terms.
N - Needs its neighbours. Any other passage: introduction, background, related work, methods, derivations, equations, experimental setup, results, discussion, figure or table captions, bullet lists, fragments.

Reply with the letter only.
\end{lstlisting}

\begin{table}[htbp]
\caption{Contextual generation prompt templates and dynamic routing decision criteria.}\label{tab:prompts}
\begin{tabular}{@{}p{3.6cm}p{9.2cm}@{}}
\toprule
Component & Specification \\
\midrule
Summariser model & \texttt{gemma3:4b-it-qat} served by Ollama; options \texttt{temperature 0}, \texttt{num\_ctx 16384} (the \texttt{SYSTEM} block and \texttt{PARAMETER} lines of \texttt{Models/chunker\_full\_doc.Modelfile}, passed per call so that the base model stays resident) \\
Summariser system prompt & Listing~\ref{lst:summariser} (789 characters) \\
Summariser messages & user~1: ``\texttt{FULL DOCUMENT:}'' + newline + the chunks of $W_i(k)$ joined by single spaces; user~2: ``\texttt{CHUNK:}'' + newline + $c_i$ \\
Embedded text & $s_i$ + blank line + $c_i$ for $k_i>0$; $c_i$ alone for $k_i=0$ (routed) \\
Router model & the same resident runner and \texttt{num\_ctx}; options \texttt{temperature 0}, \texttt{num\_predict 1}, \texttt{top\_p 1}; \texttt{logprobs} with \texttt{top\_logprobs 20} \\
Router system prompt & Listing~\ref{lst:router} \\
Router user message & ten blocks ``\texttt{Passage:} $x$ / \texttt{Answer:} S$\mid$N'' followed by ``\texttt{Passage:} $\mathrm{clip}(c_i)$ / \texttt{Answer:}'' \\
Query clipping & first 1\,500 and last 500 characters, joined by ``\texttt{ [...] }'' when the chunk exceeds 2\,020 characters \\
Exemplars (held-out papers) & abstract (S); formula block with ``where \ldots'' (N); related work (N); dangling caption fragment (N); reference-list excerpt (S); results pointing to a figure (N); device description (N); conclusion (S); experimental setup (N); acknowledgement (S) \\
Decision & argmax of the letter log-probabilities; $\kappa(\mathrm S)=0$ (no summariser call), $\kappa(\mathrm N)=2$; fallback: parse the emitted letter, else $k_{\mathrm{default}}=2$ \\
Tier-1 criteria & $\ge 3$ cells ``\texttt{, \textless col\textgreater = \textless val\textgreater.}'' on one line with $\le 100$ characters per ``\texttt{ = }'', or a caption ``\texttt{TABLE I}'' / ``\texttt{Table 1:}'' / ``\texttt{Table 1 |}'' at the start of a line $\Rightarrow k_T=2$ \\
Ablated variant & four-class genre prompt (A: whole-paper summary or boilerplate $\to 0$, B: self-contained exposition $\to 1$, C: section-internal technical body $\to 2$, D: fragment $\to 3$), tables $\to 3$ in the first version and $\to 1$ in the confirmation runs (Sect.~\ref{sec:routing-results}) \\
\botrule
\end{tabular}
\footnotetext{\verify{Every row transcribes constants of \texttt{RnD/utils/dynamic\_slice\_prediction.py} (\texttt{OLLAMA\_OPTIONS}, \texttt{TOP\_LOGPROBS}, \texttt{QUERY\_HEAD/TAIL}, \texttt{LETTER\_TO\_K}, \texttt{TABLE\_K}, \texttt{\_CELL\_RE}, \texttt{\_CAPTION\_RE}, \texttt{default\_few\_shots()}) and of cells~1--2 of \texttt{RnD/ablation\_dynamic\_doc\_slice\_generation.ipynb} (\texttt{load\_modelfile}, message construction, \texttt{SKIP\_CONTEXT\_FOR\_K0}).}}
\end{table}

\subsection{Context-enriched embedding and hybrid retrieval}\label{sec:retrieval}

The enriched text of every chunk is embedded with \texttt{nomic-embed-text-v1.5} \cite{nussbaum2024nomic}, a 768-dimensional open-weight encoder chosen for its long native context (the enriched chunks reach 2\,000 tokens plus the summary), its competitive MTEB \cite{muennighoff2023mteb} score at a size that runs unquantised on the same GPU, and its permissive licence. The encoder is executed in PyTorch through LanceDB's embedding registry \cite{lancedb} on CUDA in 32-bit precision with mean pooling; the implementation applies neither the model's task prefixes nor vector normalisation, and identical settings are used for every condition, so any absolute effect of these choices cancels in the comparisons.\verify{registry call \texttt{get\_registry().get("huggingface").create(name=..., trust\_remote\_code=True, device="cuda")} in cell~5 of \texttt{RnD/ablation\_dynamic\_doc\_slice\_generation.ipynb}; pooling and the absence of prefixes/normalisation in \texttt{lancedb/embeddings/transformers.py} of the installed LanceDB 0.26.1; the vector width 768 is stored in every \texttt{RnD/db/*.lance} manifest} The LanceDB table holds the enriched text, its vector, the raw chunk, the summary, the document name, the content hash and, for the routed pipeline, the radius and the tier that produced it; a BM25 full-text index is built on the enriched text and a scalar index on the hash.

Retrieval is hybrid \cite{kiran2025hybrid,wang2024searching}: for a question $q$ the vector leg returns the 20 nearest chunks by exact L2 search over the embeddings and the lexical leg the 20 best BM25 matches \cite{robertson1995okapi,robertson2009bm25}; the union is re-scored by a ColBERTv2 late-interaction re-ranker \cite{khattab2020colbert,santhanam2022colbertv2} (\texttt{colbert-ir/colbertv2.0}) on the enriched text and cut to the top 20. Every system in this article, including the full-document baseline and the no-context control, is indexed and queried with exactly this procedure, so the only experimental variable is the text that was embedded. Because the GPU cannot host the Ollama runner and the PyTorch models at the same time without spilling, the pipeline runs the phases strictly in sequence (routing and summarisation, then \texttt{ollama stop}, then embedding and indexing, then retrieval); Table~\ref{tab:hardware} lists the execution environment.

\begin{table}[htbp]
\caption{Experimental hardware and software environment and embedding configuration.}\label{tab:hardware}
\begin{tabular}{@{}p{4.2cm}p{8.6cm}@{}}
\toprule
Item & Specification \\
\midrule
CPU & Intel Core i7-14650HX, 24 threads \\
GPU & NVIDIA GeForce RTX 4060 Laptop GPU, 8\,GB VRAM; memory bandwidth \todo{[insert nominal GB/s]} \\
System memory & 31.06\,GB \\
Operating system / Python & Linux (zen kernels 6.18--7.2 across the measurement dates); Python 3.13 (2026 spring and summer runs) and 3.12 (router study) \\
Inference server & Ollama \todo{[insert server version, e.g. from \texttt{ollama --version}]}; one resident runner during generation \\
Summariser and router model & \texttt{gemma3:4b-it-qat} (Gemma~3, 4B parameters, quantisation-aware trained int4 weights \cite{gemma3report2025,jacob2018quantization}); layer count, KV heads and head width \todo{[verify from the model config: $n_L$, $n_{kv}$, $d_h$, local/global attention pattern]} \\
Context window (\texttt{num\_ctx}) & 16\,384 tokens for the slice and routed pipelines; 30\,000 tokens for the full-document baseline (\texttt{Models/chunker\_anthropic.Modelfile}) \\
Decoding & temperature 0; router additionally \texttt{num\_predict 1} with top-20 log-probabilities \\
Embedding model & \texttt{nomic-ai/nomic-embed-text-v1.5}, 768-d, mean pooling, fp32 on CUDA via the LanceDB HuggingFace registry \\
Vector store / search & LanceDB 0.26.1; hybrid = exact-L2 vector top-20 $\cup$ BM25 top-20 $\rightarrow$ ColBERTv2 re-rank $\rightarrow$ 20 \\
Parsing / chunking & docling 2.68.0, docling-core 2.59.0 (HybridChunker, RapidOCR); transformers 4.57.5; torch 2.9.1 \\
Energy tracking & CodeCarbon 3.2.3 (spring/summer 2026 rows) and 3.3.1 (September 2026 rows) \\
Peak GPU memory & 7\,695\,MiB over the full pipeline; generation phase $\approx 4.6$\,GB for the runner plus $\approx 1$\,GB for Docling when run concurrently \\
\botrule
\end{tabular}
\footnotetext{\verify{CPU, GPU, RAM, OS, Python and CodeCarbon versions: columns \texttt{cpu\_model}, \texttt{gpu\_model}, \texttt{ram\_total\_size}, \texttt{os}, \texttt{python\_version}, \texttt{codecarbon\_version} of \texttt{RnD/emissions\_data/emissions.csv} (note that rows written by CodeCarbon 3.3.1 are shifted by two columns from \texttt{region} onward relative to the 3.2.3 header; read those rows by position). Library versions: \texttt{RnD/requirements.txt}. Context windows: the two Modelfiles in \texttt{RnD/Models/}. Peak memory: \texttt{README.md} line~59 and the comment in cell~4 of \texttt{RnD/doc\_slice\_chunking.ipynb} (``4644 MB (ollama) + \ldots around 1 GB'').}}
\end{table}

\subsection{Computational complexity and memory analysis}\label{sec:complexity}

Index-time contextualisation issues one summariser call per chunk, so the cost of a document is the sum of the prompt costs of Eq.~\eqref{eq:prompt}. Writing $\bar\ell=L/n$ for the mean chunk length and neglecting the fixed system prompt $s$, the prompt-token volumes of the three designs are
\begin{equation}\label{eq:volume}
C_{\mathrm{full}}=\sum_{i=1}^{n}(L+\ell_i)=(n+1)L,\qquad
C_{k}=\sum_{i=1}^{n}\Bigl(\sum_{j\in W_i(k)}\ell_j+\ell_i\Bigr)\le n\,(2k+2)\,\ell_{\max},\qquad
C_{\mathrm{dyn}}=\sum_{i:\,k_i>0}\Bigl(\sum_{j\in W_i(k_i)}\ell_j+\ell_i\Bigr)+\sum_{i\notin\mathcal T}\rho_i,
\end{equation}
where $\mathcal T$ is the set of table-gated chunks and $\rho_i$ the router prompt of chunk $i$ (about 2\,000 tokens, of which the exemplar prefix is shared and served from the cache).\verify{router prompt length: mean 2\,062 tokens printed by \texttt{RnD/ollama\_dynamic\_slice\_length.ipynb}} The full-document volume is quadratic in the number of chunks, $(n+1)L=\Theta(n^2\bar\ell)$, whereas the bounded volume is linear, $C_k=\Theta(n\bar\ell)$, with a constant $2k+2$ that the router lowers further by skipping chunks.

Attention makes the gap wider. For a decoder-only transformer with $N_{\mathrm{par}}$ non-embedding parameters, $n_L$ layers and attention width $d_{\mathrm{attn}}$, the forward pass over a prompt of $P$ tokens costs approximately \cite{kaplan2020scaling}
\begin{equation}\label{eq:flops}
F(P)\approx 2N_{\mathrm{par}}\,P+2\,n_L\,d_{\mathrm{attn}}\,P^{2},
\end{equation}
so the attention term of a document is $\propto\sum_i P_i^2$.

\begin{proposition}[Asymptotic index-time cost]\label{prop:complexity}
Let every chunk satisfy $\ell_j\le\ell_{\max}$ and let $n\ge 2k+1$. Then the attention work of full-document contextualisation grows as $\sum_i (L+\ell_i)^2=\Theta(n^{3}\bar\ell^{\,2})$, while that of Document Slice contextualisation with radius $k$ is bounded by $\sum_i P_i(k)^2\le n\,(2k+2)^2\ell_{\max}^2=\Theta(n\,\ell_{\max}^2)$, independent of the document length beyond $n$. The two designs differ by a factor $\Theta\bigl(n^2/(2k+2)^2\bigr)$ in attention work and $\Theta\bigl(n/(2k+2)\bigr)$ in prompt tokens.
\end{proposition}
\begin{proof}
For every $i$, $L\le L+\ell_i\le 2L$ and $L=n\bar\ell$, hence $nL^2\le\sum_i(L+\ell_i)^2\le 4nL^2$, i.e.\ $\sum_i(L+\ell_i)^2=\Theta(n^3\bar\ell^{\,2})$. By Eq.~\eqref{eq:window} the window holds $\min(2k+1,n)=2k+1$ chunks, so $P_i(k)\le(2k+1)\ell_{\max}+\ell_{\max}$ and $\sum_i P_i(k)^2\le n(2k+2)^2\ell_{\max}^2$; the lower bound $\sum_i P_i(k)^2\ge n\bar\ell^{\,2}$ follows from $P_i(k)\ge\ell_i$ and the Cauchy--Schwarz inequality, so the sum is $\Theta(n\ell_{\max}^2)$ whenever $\bar\ell$ and $\ell_{\max}$ are of the same order, as they are for token-bounded chunkers. The token statement follows from Eq.~\eqref{eq:volume}.
\end{proof}

On the benchmark corpus the bounds are far from tight but the ordering is the same (Fig.~\ref{fig:complexity}): summing Eq.~\eqref{eq:prompt} over the 382 chunks with the actual chunk texts, the full-document design sends 18.73\,M characters ($\approx 4.5$\,M tokens) to the summariser, fixed radii $k=1,2,3,4$ send 4.25, 6.36, 8.43 and 10.32\,M, and the routed pipeline sends 4.39\,M ($23.4$\,\% of the baseline and $69.0$\,\% of fixed $k=2$) in 261 instead of 382 calls; the quadratic proxy $\sum_i P_i^2$ is 4.6 times smaller for $k=3$ and 10.8 times smaller for the routed pipeline than for the baseline.\verify{re-implementation of the \texttt{cost()} formula of \texttt{RnD/dynamic\_slice\_length/router\_experiments/scripts/cost.py} (window characters of Eq.~\eqref{eq:window} plus the chunk, summed over \texttt{RnD/split\_documents/*.json}; the full-document prompt is the space-joined document plus the chunk; the routed radii are the \texttt{doc\_slice\_radius} fields of \texttt{RnD/preprocessed\_chunks/ablation\_doc\_slice\_radius\_dynamic.json}): 18\,728\,428 / 4\,254\,674 / 6\,360\,511 / 8\,433\,157 / 10\,320\,971 / 4\,391\,194 characters; $\sum P_i^2$ ratios 4.56 and 10.81; tokens $\approx$ characters$/4.17$} The bound of Eq.~\eqref{eq:prompt} is rarely approached: the mean $k=2$ prompt over the corpus is about 4\,200 tokens and a window of median-length chunks about 2\,900, against the 12\,190-token bound, because the chunker fills its 2\,000-token budget for only a minority of chunks (Table~\ref{tab:parsing}).\verify{$6\,360\,511/382/4.17+190\approx 4\,180$ tokens (fixed-$k=2$ prompt volume of the \texttt{cost.py} re-implementation divided by 382 calls and by 4.17 characters per token, plus the 190-token system prompt); $6\times450+190=2\,890$ with the token median of Table~\ref{tab:parsing}; script: \texttt{RnD/verification/prompt\_volume\_and\_attention\_proxy.py}} The measured generation-time ratio between the baseline and the routed pipeline, 6.3 (Sect.~\ref{sec:primary}), lies between the linear proxy (4.3) and the quadratic proxy (10.8), as expected for prompts of a few thousand tokens where the linear term of Eq.~\eqref{eq:flops} still matters and where the baseline's paging (below) adds a cost that neither proxy models.

\begin{figure}[htbp]
  \centering
  \todo{[FIGURE INSERTION: two-panel line chart. Panel (a): index-time prompt volume per document as a function of the number of chunks $n$ (x-axis 8--33, the range of the corpus), series: full document $(n+1)L$ with $L=n\bar\ell$, fixed $k=1,2,3$ from Eq.~\eqref{eq:volume}, and the routed pipeline with the measured skip rate (32\,\% of chunks, $k=2$ otherwise); y-axis in tokens (use $\bar\ell=664$). Panel (b): the quadratic attention proxy $\sum_i P_i^2$ for the same series, log y-axis. Overlay the corpus-level measured points: 4.49\,M / 1.02\,M / 1.53\,M / 2.02\,M / 1.05\,M tokens (full, $k=1$, $k=2$, $k=3$, routed) and the ratios 4.6$\times$ and 10.8$\times$. Data source: the \texttt{cost.py} re-implementation described in the text.]}
  \caption{Theoretical index-time cost curves. Full-document contextualisation grows quadratically in the number of chunks in prompt tokens and cubically in attention work, whereas bounded slices grow linearly with a constant set by the radius; the routing agent lowers that constant further by leaving self-describing chunks un-summarised. Over the 382-chunk corpus the routed pipeline sends 23.4\,\% of the baseline's prompt text and incurs about one tenth of its quadratic attention proxy; the measured 6.3-fold reduction of generation time lies between the two proxies.\verify{corpus volumes and ratios as in the preceding paragraph; 6.3 $= 3294.3/525.3$ from the \texttt{duration} column of the rows timestamped 2026-03-05T23:59:43 and 2026-09-25T20:06:47 in \texttt{RnD/emissions\_data/emissions.csv}}}
  \label{fig:complexity}
\end{figure}

Memory decides whether the arithmetic can be done on the device at all. During generation the GPU must hold the quantised weights $M_w$, the activations $M_{\mathrm{act}}$ and the key--value (KV) cache, which for a context of $P$ tokens occupies
\begin{equation}\label{eq:kv}
M_{\mathrm{KV}}(P)=2\,n_L\,n_{kv}\,d_h\,b\,P
\end{equation}
bytes ($n_{kv}$ key--value heads of width $d_h$, $b$ bytes per element, the factor 2 for keys and values) \cite{pope2023efficiently,ainslie2023gqa}. The run stays on the device only if the cache of Eq.~\eqref{eq:kv} fits next to the weights,
\begin{equation}\label{eq:vram}
M_w+M_{\mathrm{KV}}(P_{\mathrm{ctx}})+M_{\mathrm{act}}\le V_{\mathrm{GPU}}
\quad\Longleftrightarrow\quad
P_{\mathrm{ctx}}\le P^{\star}=\frac{V_{\mathrm{GPU}}-M_w-M_{\mathrm{act}}}{2\,n_L\,n_{kv}\,d_h\,b},
\end{equation}
where $P_{\mathrm{ctx}}$ is the context window the runtime reserves, not the length of the current prompt: Ollama allocates the cache for the configured \texttt{num\_ctx} when the model is loaded. With the bounded slice, $P_i(k)\le s+(2k+2)\ell_{\max}$ (Eq.~\eqref{eq:prompt}) is independent of the document, so a window of 16\,384 tokens covers every prompt of the corpus for $k\le3$ and the requirement of Eq.~\eqref{eq:vram} is fixed once and for all. The full-document design needs $P_{\mathrm{ctx}}\ge s+L_{\max}+\ell_{\max}$, which for the longest paper of the corpus (about 18\,300 tokens) already exceeds 16\,384 and forced a 30\,000-token window.\verify{longest document 76\,161 characters (sum of chunk lengths of \texttt{s41586-019-1138-y.pdf.json} in \texttt{RnD/split\_documents/}), divided by 4.17 characters per token; windows from \texttt{RnD/Models/chunker\_anthropic.Modelfile} (30000) and \texttt{chunker\_full\_doc.Modelfile} (16384)} On the 8\,GB device this crossed $P^{\star}$: the runtime fell back to unified memory and streamed part of the working set from host RAM over PCIe on every decoding step (Fig.~\ref{fig:memory}). In the conference study we observed host-memory growth of roughly 10\,GB and a resident footprint of about 5\,GB for the slice pipeline; those observations were not logged and are marked for re-measurement.\todo{[unlogged observations from \texttt{citds.tex}, Sect.~``Hardware Profiling''; no artefact in the repository]} What the energy tracker did log is consistent with a bandwidth-bound baseline: its mean GPU utilisation was 46.5\,\% over the run, against 60.6\,\% for the fixed $k=3$ run and 90.1\,\% for the routed run.\verify{column \texttt{gpu\_utilization\_percent} of the rows timestamped 2026-03-05T23:59:43 (46.5), 2026-03-05T22:53:27 (60.6) and 2026-09-25T20:06:47 (90.1) in \texttt{RnD/emissions\_data/emissions.csv}; whole-machine averages sampled every second by CodeCarbon (\texttt{tracking\_mode = machine})} \todo{[VERIFY the paging mechanism from the Ollama server log of a baseline run (``offloaded $x/y$ layers to GPU'' and the reported KV-cache size), and re-measure resident VRAM and host RSS with nvidia-smi/psutil for the 30\,000- and 16\,384-token windows; Table~\ref{tab:hardware} needs $n_L$, $n_{kv}$, $d_h$ of Gemma~3~4B to evaluate Eq.~\eqref{eq:vram} numerically.]}

\begin{figure}[htbp]
  \centering
  \todo{[FIGURE INSERTION: memory-topology diagram. Left: GPU with 8\,GB VRAM containing the int4 weights ($M_w$), activations and the KV cache sized by \texttt{num\_ctx}; right: 31\,GB host RAM; between them the PCIe link. Show two stacked-bar insets: (i) slice/routed pipeline with \texttt{num\_ctx} = 16\,384 fitting under the 8\,GB line; (ii) full-document baseline with \texttt{num\_ctx} = 30\,000 crossing the line, with the overflow mapped to host memory and an arrow labelled ``per-step PCIe traffic''. Annotate the threshold $P^{\star}$ of Eq.~\eqref{eq:vram}. Values for the bar heights are \todo{[to be measured]}; the qualitative layout follows the conference paper.]}
  \caption{Unified-memory offloading model. When the reserved context window pushes weights plus KV cache beyond the device memory, the runtime maps the overflow to host memory and every decoding step crosses the PCIe bus, which shows up as low GPU utilisation and long wall-clock times rather than as an out-of-memory failure. The bounded slice keeps the reserved window, and therefore the memory footprint, independent of document length.}
  \label{fig:memory}
\end{figure}

\subsection{Evaluation metrics and statistical protocol}\label{sec:metrics}

For a question $q\in Q$ let $G_q$ be its set of gold supporting chunks, $R_q^K$ the ordered list of the top-$K$ retrieved chunks and $\rnk_q(c)$ the 1-based position of chunk $c$ in that list. We report
\begin{align}
\Rec{K}&=\frac{1}{|Q|}\sum_{q\in Q}\frac{|R_q^K\cap G_q|}{|G_q|},\qquad
Q_K^{+}=\{\,q\in Q : R_q^K\cap G_q\neq\varnothing\,\},\label{eq:recall}\\
\MRR{K}&=\frac{1}{|Q_K^{+}|}\sum_{q\in Q_K^{+}}\frac{1}{\min_{c\in R_q^K\cap G_q}\rnk_q(c)},\qquad
\MAR{K}=\frac{1}{|Q_K^{+}|}\sum_{q\in Q_K^{+}}\frac{1}{|R_q^K\cap G_q|}\sum_{c\in R_q^K\cap G_q}\rnk_q(c),\label{eq:mrr}
\end{align}
for $K\in\{5,10,15,20\}$. Recall is the primary metric because every question requires at least two chunks and a generator can only use what was retrieved. The mean reciprocal rank \cite{voorhees1999trec8} reflects the position of the first relevant chunk, and the mean average rank (MAR) is our exploratory measure of attentional accessibility: it is the mean position of \emph{all} retrieved gold chunks and tells whether the supporting evidence sits in the top part of the context that a generator will read, where the positional bias of Sect.~\ref{sec:intro} is least harmful. Two implementation details matter for comparability. First, the ranking metrics are averaged over $Q_K^{+}$, i.e.\ questions with no gold chunk in the top $K$ are excluded rather than scored zero, exactly as in the evaluation notebooks; we therefore report $|Q_K^{+}|$ alongside them (for the routed pipeline 233, 246, 247 and 248 of the 250 questions at $K=5,10,15,20$).\verify{``Number of N values with NaN MRR in own metrics: 17/250, 4/250, 3/250, 2/250'' printed by cell~7 of \texttt{RnD/benchmarking\_retrievals.ipynb}; the \texttt{np.nanmean} calls in cell~9 implement the exclusion} Second, gold sets are compared by content hash, so a chunk retrieved from any condition's table is matched to the question set without string comparison.

Differences between systems are tested per question. For \Rec{20} we compute the paired difference of means with a bias-corrected and accelerated (BCa) bootstrap confidence interval at the 99\,\% level \cite{efron1987better,smucker2007comparison} (SciPy \cite{virtanen2020scipy}, 9\,999 resamples, fixed seed), and we apply an exact McNemar test \cite{McNemar_1947,dietterich1998approximate} to the binary event ``all gold chunks of the question are in the top 20'', whose discordant counts $b$ (only the compared system succeeds) and $c$ (only the baseline succeeds) are reported with the $p$-value. Because the summariser is a stochastic component even at temperature~0 (its outputs depend on the state of the server's cache and batching), we also estimate the noise floor of a single run by regenerating one configuration and counting the questions whose top-20 list changes; the result (Sect.~\ref{sec:routing-results}) is that a single run resolves differences of roughly 0.7 points of \Rec{20} and 2 points of \MRR{20}, and we treat smaller gaps as ties.

\subsection{Datasets}\label{sec:datasets}

\subsubsection{Scientific benchmark}
The primary corpus consists of 25 open-access papers from IEEE and Nature-portfolio journals across engineering, physics, computing and the life sciences, chosen for their heavy use of tables, formulas and figure references, i.e.\ for the layout artefacts that make chunking hard. Following the LLM-as-a-judge paradigm \cite{zheng2023judgingllmasajudgemtbenchchatbot} we built a ``silver standard'' question set \cite{rebholz2010silverstandard}: Gemini~3~Pro was given the chunk JSON of one paper at a time and instructed, in the role of a senior NLP data scientist, to write exactly ten high-complexity question--answer pairs grounded only in the provided text, each requiring at least two distinct chunk identifiers, with a plain-text gold answer and the list of supporting chunk identifiers (the abridged prompt is reproduced in the conference paper \cite{sandor2025documentslicecontextualretrieval}). A validator script rejects questions with fewer than two supporting chunks or with identifiers that do not exist in the paper's chunk file. The result is 250 questions with 530 gold references over 226 distinct chunks (Table~\ref{tab:corpora}). Because the gold labels were produced by a large proprietary model while the summaries are written by a 4-billion-parameter open model of different provenance, a shared preference for particular chunks is unlikely to be a model-level artefact; the question set nonetheless remains a silver rather than a gold standard.

\subsubsection{Policy corpus}
To test the method outside scientific prose and inside the target domain of regulatory text, we use four public European Commission documents on tobacco-control policy: two Commission reports to the Parliament and the Council (COM(2016)~269 on refillable electronic cigarettes and COM(2021)~249 on the application of Directive 2014/40/EU), the implementation report on the Council Recommendation on smoke-free environments (SWD(2013)~56) and the methodology report of the Independent Advisory Panel on characterising flavours (2023). They were chunked with the same recipe and questioned with the same protocol (80 questions, 164 gold references); a random 20\,\% subset of the questions was checked manually for semantic alignment with the supporting chunks \todo{[confirm: the manual validation of the policy questions is described in the conference paper but no record of it exists in the repository]}. Although the subject matter is health rather than energy policy, the documents share the genre that energy-governance workflows deal with: directive articles, implementation tables by member state, methodology sections and annexes.

\subsubsection{Held-out exemplar papers}
The ten few-shot exemplars of the router come from six further papers that appear in neither benchmark (109 chunks), so that no text the router has seen as an example can be retrieved as an answer.\verify{\texttt{RnD/split\_documents\_router/}: 6 files, 109 records; \texttt{RnD/input\_router/} holds the PDFs}

\begin{table}[htbp]
\caption{Evaluation corpora: the scientific silver-standard benchmark and the European Commission policy set.}\label{tab:corpora}
\begin{tabular}{@{}lll@{}}
\toprule
Quantity & Scientific benchmark & EC policy set \\
\midrule
Documents & 25 (IEEE, Nature portfolio) & 4 (COM(2016)~269, COM(2021)~249, SWD(2013)~56, IAP methodology 2023) \\
Chunks & 382 & 141 (14 / 32 / 35 / 60 per document) \\
Characters (total) & 1\,056\,773 & 259\,985 \\
Median chunk length (characters) & 2\,090 & 1\,206 \\
Chunk length, tokens (mean / median) & 663.6 / 450 & 446.8 / 246 \\
Questions & 250 (10 per document) & 80 (10 / 20 / 20 / 30 per document) \\
Supporting chunks per question & 2--4 (mean 2.12; 223 / 24 / 3 questions with 2 / 3 / 4) & 2--3 (mean 2.05; 76 / 4) \\
Gold references & 530 & 164 \\
Distinct gold chunks (share of corpus) & 226 (59.5\,\%) & 94 (66.7\,\%) \\
Question generator & Gemini~3~Pro & Gemini~3~Pro \\
Manual validation & --- & 20\,\% random subset \todo{[confirm]} \\
Router exemplar pool & 6 held-out papers, 109 chunks & --- \\
\botrule
\end{tabular}
\footnotetext{\verify{Chunk counts, character statistics and per-document chunk counts: read-only pass over \texttt{RnD/split\_documents/*.json} and \texttt{RnD/split\_documents/smoking/*.json}; question statistics: \texttt{RnD/q\_and\_a/Gemini/scientific\_multi\_chunk\_control.json} and \texttt{RnD/q\_and\_a/Gemini/smoking.json} (count the \texttt{supporting\_chunks} lists; distinct gold ids as a set); token statistics: \texttt{len(tokenizer.tokenize(text))} with the nomic tokenizer, as for Table~\ref{tab:parsing}, over \texttt{RnD/split\_documents/*.json} (mean 663.58, median 450, total 253\,488) and over the final 141 records of \texttt{RnD/split\_documents/smoking/*.json} (mean 446.80, median 246, min 1, max 1\,992, total 62\,999); the 420.94 printed by cell~2 of \texttt{RnD/doc\_slice\_chunking.ipynb} refers to an earlier 144-chunk version of the policy corpus; scripts: \texttt{RnD/verification/chunk\_character\_statistics.py} and \texttt{chunk\_token\_statistics.py}; exemplar pool: \texttt{RnD/split\_documents\_router/} (6 files, 109 records).}}
\end{table}

\subsection{Continuous energy and carbon accounting}\label{sec:energy}

Every generation run is wrapped in a CodeCarbon tracker \cite{codecarbon_software,lacoste2019quantifyingcarbonemissionsmachine} that samples the instantaneous power of the GPU, the CPU and the memory once per second and integrates it over the run (Eq.~\eqref{eq:energy}); emissions follow from the energy through the grid factor (Eq.~\eqref{eq:emissions}):
\begin{equation}\label{eq:energy}
E_{\mathrm{total}}=\int_{0}^{T}\bigl(P_{\mathrm{GPU}}(t)+P_{\mathrm{CPU}}(t)+P_{\mathrm{DRAM}}(t)\bigr)\,dt
\;\approx\;\sum_{t=1}^{\lceil T/\Delta t\rceil}\bigl(P^{(t)}_{\mathrm{GPU}}+P^{(t)}_{\mathrm{CPU}}+P^{(t)}_{\mathrm{DRAM}}\bigr)\,\Delta t,\qquad \Delta t=1\,\mathrm{s},
\end{equation}
\begin{equation}\label{eq:emissions}
\mathrm{Emissions}=E_{\mathrm{total}}\times\mathrm{CI}_{\mathrm{grid}}\times\mathrm{PUE},
\end{equation}
where $\mathrm{CI}_{\mathrm{grid}}$ is the operational carbon intensity of the regional grid in \gco{} per kWh and the power-usage effectiveness is 1 for a workstation. GPU power is read from the NVIDIA management library; CPU power was read from the processor's RAPL counters for the spring-2026 runs and, after a permission change of the operating system, estimated from CPU load and the processor's thermal design power for the later runs; memory power is CodeCarbon's constant estimate. The grid factor for Hungary in the tracker's data set is 204.19\,\gco{}/kWh for the 3.2.3 releases; the 3.3.1 releases used for the router study yield 211.5--214.0\,\gco{}/kWh on the same machine. Table~\ref{tab:power} lists the coefficients. Two consequences follow for the comparisons of Sect.~\ref{sec:results}: energy is the primary efficiency quantity and emissions are derived from it, and because the CPU measurement method changed between the baseline runs and the router runs, the GPU energy, which was measured identically throughout, is reported next to the total as the component that is comparable without qualification.\verify{measurement methods: the CodeCarbon start-up messages saved in \texttt{RnD/anthropic\_traditional\_chunking.ipynb} (March: ``Tracking Intel CPU via RAPL interface''; July: ``RAPL -- Permission denied \ldots Falling back on CPU load mode''); the grid factor is \texttt{emissions}/\texttt{energy\_consumed} of each row of \texttt{RnD/emissions\_data/emissions.csv} (204.19 for the four 3.2.3 rows; 211.50--214.03 for the complete 3.3.1 rows)}

Durations are reported in two scopes that must not be mixed. The \emph{tracked generation scope} is the interval wrapped by the energy tracker: the routing and summarisation loop over the pre-chunked documents and the final unloading of the model, without PDF parsing, embedding or indexing; all energy and emission figures refer to this scope. The \emph{end-to-end preprocessing scope} is the wall-clock time of the complete notebook cell that also converts and chunks the PDFs with Docling; the conference paper's headline times (71\,min~34\,s for the baseline, 29\,min~09\,s for the fixed $k=3$ pipeline) belong to this scope and were not energy-tracked. \todo{[confirm the tracked scope of the fixed $k=3$ energy row of 5 March 2026: the notebook version of that date (commit fe18290 of \texttt{doc\_slice\_chunking.ipynb}) appears to run the Docling conversion inside the tracked block, whereas the baseline row wraps only the summariser loop; if so, the static $k=3$ energy in Tables~\ref{tab:sweep} and \ref{tab:primary} is an upper bound.]}

\begin{table}[htbp]
\caption{Power-measurement coefficients, per-run power draw and grid emission factors.}\label{tab:power}
\begin{tabular}{@{}p{4.6cm}p{8.2cm}@{}}
\toprule
Quantity & Value \\
\midrule
Sampling interval / tracking mode / PUE & 1\,s / whole machine / 1.0 \\
GPU power & measured via NVML (all runs) \\
CPU power & RAPL counters (runs of 5 March 2026); load-based estimate with the tracker's TDP fallback (``4\,W per thread $\times$ 24 threads $=$ 96\,W'') from 19 July 2026 on \\
Memory power & constant 20.0\,W estimate (31\,GB installed) \\
Grid carbon intensity (Hungary) & 204.19\,\gco{}/kWh (CodeCarbon 3.2.3 data set); effective 211.5--214.0\,\gco{}/kWh in the 3.3.1 rows \\
Mean GPU power / utilisation: full-document baseline & 51.5\,W / 46.5\,\% \\
Mean GPU power / utilisation: static $k=3$ & 55.0\,W / 60.6\,\% \\
Mean GPU power / utilisation: routed pipeline & 62.5\,W / 90.1\,\% \\
Mean CPU power: baseline / $k=3$ (RAPL) & 46.1\,W / 58.0\,W \\
Mean CPU power: routed pipeline (estimate) & 10.4\,W \\
Mean host memory in use: baseline / $k=3$ / routed & 15.9 / 19.4 / 17.3\,GB \\
Water footprint & not tracked (WUE 0) \todo{[add if the journal requires it]} \\
\botrule
\end{tabular}
\footnotetext{\verify{All rows are columns of \texttt{RnD/emissions\_data/emissions.csv}: \texttt{gpu\_power}, \texttt{gpu\_utilization\_percent}, \texttt{cpu\_power}, \texttt{ram\_used\_gb}, \texttt{ram\_power}, \texttt{pue}, \texttt{tracking\_mode} of the rows timestamped 2026-03-05T23:59:43 (baseline), 2026-03-05T22:53:27 (static $k=3$) and 2026-09-25T20:06:47 (routed); the TDP fallback message is in the saved CodeCarbon log of \texttt{RnD/anthropic\_traditional\_chunking.ipynb}. Rows written by CodeCarbon 3.3.1 must be read by position (two-column shift after \texttt{region}).}}
\end{table}

\section{Results and empirical evaluation}\label{sec:results}

All metrics are those of Sect.~\ref{sec:metrics} (Eqs.~\eqref{eq:recall} and \eqref{eq:mrr}) on the corpora of Sect.~\ref{sec:datasets}. Section~\ref{sec:static} reports the static radius sweep, Sect.~\ref{sec:primary} the primary comparison, Sect.~\ref{sec:routing-results} the routing statistics and the noise floor, Sect.~\ref{sec:ranking} the ranking dynamics, Sect.~\ref{sec:policy} the policy corpus and Sect.~\ref{sec:profiling} the hardware profile.

\subsection{Static radius sweep}\label{sec:static}

Table~\ref{tab:sweep} and Fig.~\ref{fig:heatmap} report the fixed-radius ablation that motivated the router. Generation time grows linearly with the radius, by roughly five minutes per unit of $k$ over the 382 chunks (7\,min~28\,s at $k=0$ to 28\,min~11\,s at $k=4$), while recall saturates: $k=2$ and $k=3$ are the best fixed radii (\Rec{20} of 0.94667 and 0.94533) and $k=4$ falls back to the level of $k=1$ (0.94333) although it sends 62\,\% more prompt text than $k=2$.\verify{times: the ``k=0 \ldots k=4'' comment block in cell~1 of \texttt{RnD/ablation\_doc\_slice\_generation.ipynb} (wall-clock of the generation loop, no tracker); recall values: ``METRICS AT N=20'' blocks of \texttt{RnD/ablation\_doc\_slice\_retrieval.ipynb} for the tables \texttt{ablation\_doc\_slice\_radius\_0..4}; prompt-text ratio $10\,320\,971/6\,360\,511=1.62$ from the \texttt{cost.py} re-implementation} The full-document baseline remains the most accurate configuration (0.95400), 0.7 points above the best fixed radius, at a tracked generation time of 54\,min~54\,s and 106.7\,Wh. Even the no-context hybrid control reaches 0.93933, so the entire effect of contextualisation on this benchmark is confined to about 1.5 points of \Rec{20}\verify{Table~\ref{tab:sweep}: $0.95400-0.93933=0.0147$}; what the ablation shows is that most of it is obtained from the immediate neighbourhood of a chunk and none from the far context, as the positional-bias literature predicts. The ranking metrics do not favour wide windows either: the best \MRR{20} (0.79900) and the best \MAR{20} (3.72605) among the fixed radii belong to $k=1$, and $k=3$ has the worst MAR of the sweep (4.01999).\verify{\MRR{20} and \MAR{20} columns of Table~\ref{tab:sweep}, from the same notebook printouts} All recall differences among $k\in\{1,\dots,4\}$ lie inside the single-run noise band established in Sect.~\ref{sec:routing-results}, so the sweep identifies a plateau rather than an optimum, which is the situation in which a per-chunk decision can save cost without losing accuracy.

\begin{table}[htbp]
\caption{Ablation over static window radii versus full-document context and the routed pipeline (250 questions, tracked generation scope).}\label{tab:sweep}
\begin{tabular}{@{}lcccccc@{}}
\toprule
Method & Generation time & Prompt volume (M chars) & Energy (Wh) & \Rec{20} & \MRR{20} & \MAR{20} \\
\midrule
Hybrid retrieval, no context & --- & 0 & 0 & 0.93933 & 0.77484 & 4.03629 \\
Document slice $k=0$ (chunk-only summary) & 7\,min~28\,s\footnotemark[1] & 2.11 & \todo{[not tracked]} & 0.93733 & 0.78612 & 3.80827 \\
Document slice $k=1$ & 12\,min~08\,s\footnotemark[1] & 4.25 & \todo{[not tracked]} & 0.94333 & 0.79900 & 3.72605 \\
Document slice $k=2$ & 17\,min~28\,s\footnotemark[1] & 6.36 & \todo{[not tracked]} & 0.94667 & 0.76799 & 3.80948 \\
Document slice $k=3$ & 22\,min~41\,s\footnotemark[1] & 8.43 & 52.52\footnotemark[2] & 0.94533 & 0.78774 & 4.01999 \\
Document slice $k=4$ & 28\,min~11\,s\footnotemark[1] & 10.32 & \todo{[not tracked]} & 0.94333 & 0.78335\footnotemark[3] & 3.95884 \\
Full-document context (baseline) & 54\,min~54\,s & 18.73 & 106.68 & 0.95400 & 0.78635 & 3.93548 \\
Routed pipeline (this work) & 8\,min~45\,s & 4.39 & 13.22 & 0.95200 & 0.79798 & 3.91062 \\
\botrule
\end{tabular}
\footnotetext[1]{Wall-clock of the generation loop as recorded in the notebook; these runs were not energy-tracked.}
\footnotetext[2]{Energy of a separate $k=3$ generation of 5 March 2026 (tracked, 23\,min~48\,s); its retrieval metrics were not evaluated, so the retrieval columns of this row come from the ablation generation. See also the scope note in Sect.~\ref{sec:energy}.}
\footnotetext[3]{The conference paper printed 0.8933 for this cell; 0.78335 is the value in the evaluation notebook.}
\footnotetext{\verify{Recall/MRR/MAR: ``METRICS AT N=20'' printouts in \texttt{RnD/ablation\_doc\_slice\_retrieval.ipynb} (tables \texttt{hybrid\_retrieval}, \texttt{ablation\_doc\_slice\_radius\_0..4}, \texttt{anthropic\_control\_table}) and in \texttt{RnD/benchmarking\_retrievals.ipynb} (routed row, own vs baseline). Times: cell~1 comments of \texttt{RnD/ablation\_doc\_slice\_generation.ipynb} for $k=0..4$; \texttt{duration} column of \texttt{RnD/emissions\_data/emissions.csv} for the baseline (3\,294.3\,s, row 2026-03-05T23:59:43) and the routed run (525.3\,s, row 2026-09-25T20:06:47). Energy: \texttt{energy\_consumed} of the same rows and of row 2026-03-05T22:53:27 ($k=3$). Prompt volumes: \texttt{cost.py} re-implementation (Sect.~\ref{sec:complexity}). Single-run noise band $\pm0.7$ points of \Rec{20}: Sect.~\ref{sec:routing-results}.}}
\end{table}

\begin{figure}[htbp]
  \centering
  \todo{[FIGURE INSERTION: heat map with columns $k\in\{0,1,2,3,4\}$ and two rows. Row ``generation time (min)'': 7.47, 12.13, 17.47, 22.68, 28.18 (from Table~\ref{tab:sweep}); row ``peak VRAM (GB)'': \todo{[not measured; record nvidia-smi peak per k]}. Optionally a third row ``prompt volume (M chars)'': 2.11, 4.25, 6.36, 8.43, 10.32. Annotate each cell with its value; use a sequential colour map per row.]}
  \caption{Cost surface of the static radius. Generation time grows linearly in $k$ (about five minutes per unit over 382 chunks) and prompt volume grows with $2k+2$, while every radius keeps the reserved context window, and hence the memory footprint, constant; the peak-memory row is a placeholder until it is measured per radius.\verify{times as in Table~\ref{tab:sweep}; prompt volumes from the \texttt{cost.py} re-implementation}}
  \label{fig:heatmap}
\end{figure}

\subsection{Primary benchmark: full-document baseline, static slice and routed pipeline}\label{sec:primary}

Table~\ref{tab:primary} is the central result. On the 250 multi-hop questions the routed pipeline reaches \Rec{20} of 0.952 against 0.954 for the full-document baseline and 0.94533 for the static $k=3$ slice, i.e.\ it is 0.2 points below the baseline and 0.7 points above the static slice, both differences inside the single-run band; at the same time it has the best \MRR{20} of the three (0.79798 against 0.78635 and 0.78774) and the best \MAR{20} (3.91062 against 3.93548 and 4.01999), and it leads both at $K=5$ (0.75167 against 0.73733 and 0.74100).\verify{routed and baseline values: the four ``METRICS AT N='' blocks of \texttt{RnD/benchmarking\_retrievals.ipynb} (OWN = \texttt{ablation\_doc\_slice\_radius\_dynamic}, BASELINE = \texttt{anthropic\_control\_table}); static $k=3$: the \texttt{ablation\_doc\_slice\_radius\_3} blocks of \texttt{RnD/ablation\_doc\_slice\_retrieval.ipynb}} The recall trajectories (Fig.~\ref{fig:recall}) show where the routed pipeline gains and loses: it retrieves more gold chunks in the first five and ten positions, falls 0.7 points behind the baseline at $K=15$ and closes to 0.2 points at $K=20$. The cost side of the table (Figs.~\ref{fig:energy} and \ref{fig:runtime}) is not inside any noise band: the routed pipeline issued 261 summariser calls instead of 382, sent 23.4\,\% of the baseline's prompt text and finished the tracked generation in 8\,min~45\,s instead of 54\,min~54\,s (6.3 times faster), consuming 13.22\,Wh instead of 106.68\,Wh ($-87.6$\,\%; GPU energy alone $-81.2$\,\%) and emitting 2.83 instead of 21.78\,\gco{} ($-87.0$\,\%).\verify{\texttt{duration}, \texttt{energy\_consumed}, \texttt{gpu\_energy}, \texttt{emissions} of the rows timestamped 2026-09-25T20:06:47 and 2026-03-05T23:59:43 in \texttt{RnD/emissions\_data/emissions.csv}; percentages $=1-\text{routed}/\text{baseline}$; calls and prompt text from the routing of \texttt{RnD/preprocessed\_chunks/ablation\_doc\_slice\_radius\_dynamic.json} (261 records with non-empty \texttt{context}) and the \texttt{cost.py} re-implementation} Against the static $k=3$ slice, which our conference paper proposed, the router removes a further 63\,\% of generation time and 75\,\% of energy (60\,\% of GPU energy) while adding 0.7 points of recall. The static slice itself already halves the baseline's energy ($-50.8$\,\%) and time ($-56.6$\,\% in the tracked scope, $-59.3$\,\% end to end) for a 0.9-point loss of \Rec{20}, which is the trade-off the conference study reported; the router turns that trade-off into one with no measurable loss.\verify{$1-52.52/106.68=0.508$, $1-1428.2/3294.3=0.566$ from the rows timestamped 2026-03-05T22:53:27 and 2026-03-05T23:59:43; end-to-end $1-1749/4294=0.593$ from the notebook comments ``71m 34s'' (\texttt{RnD/anthropic\_traditional\_chunking.ipynb}) and ``29m 9s'' (\texttt{RnD/doc\_slice\_chunking.ipynb})}

The same configuration was generated twice, independently and from a freshly loaded model: the earlier run (the replicate row of Table~\ref{tab:primary}) also reached \Rec{20} of 0.952, with \MRR{20} 0.79791 and \MAR{20} 3.92171, in 543.7\,s and 13.51\,Wh, and the two runs differ only in the third decimal of the ranking metrics. Per-question tests were computed on the replicate, whose per-question results are stored, and are reported as provisional until they are recomputed on the run in the table: \todo{[provisional, computed from \texttt{RnD/dynamic\_slice\_length/router\_experiments/results/res\_real\_binary\_t2\_heldout.json} against \texttt{res\_anthropic\_control\_table.json}: paired difference of mean \Rec{20} (baseline minus routed) $+0.0020$, 99\,\% BCa bootstrap interval $[-0.0153,\,+0.0200]$ (9\,999 resamples, seed 42, SciPy 1.17); McNemar on ``all gold chunks in the top 20'': 222 questions solved by both, 5 only by the routed pipeline, 6 only by the baseline, 17 by neither, exact $p=1.00$; per question the routed pipeline is better on 6 and worse on 7. For the static $k=3$ slice against the baseline: difference $+0.0087$, interval $[-0.0067,\,+0.0260]$, discordant counts 3 and 6, $p=0.51$. Recompute for the primary run once its per-question results are saved; script: \texttt{RnD/verification/retrieval\_significance\_tests.py}.]} Both intervals contain zero at the 99\,\% level, and the discordant counts show that the baseline and the routed pipeline fail on almost disjoint, very small sets of questions.

\begin{table}[htbp]
\caption{Primary benchmark on 250 multi-hop questions: full-document contextualisation versus the static $k=3$ slice and the routed pipeline.}\label{tab:primary}
\begin{tabular}{@{}lccccc@{}}
\toprule
Metric & Full document & Static $k=3$ & Routed pipeline & $\Delta$ routed vs full & $\Delta$ routed vs $k=3$ \\
\midrule
\Rec{5} & 0.73733 & 0.74100 & 0.75167 & $+1.4$\,pt & $+1.1$\,pt \\
\Rec{10} & 0.86500 & 0.85900 & 0.86700 & $+0.2$\,pt & $+0.8$\,pt \\
\Rec{15} & 0.93333 & 0.91067 & 0.92600 & $-0.7$\,pt & $+1.5$\,pt \\
\Rec{20} & 0.95400 & 0.94533 & 0.95200 & $-0.2$\,pt & $+0.7$\,pt \\
\MRR{20} & 0.78635 & 0.78774 & 0.79798 & $+1.5$\,\% & $+1.3$\,\% \\
\MAR{20} & 3.93548 & 4.01999 & 3.91062 & $-0.6$\,\% & $-2.7$\,\% \\
Questions with a hit in the top 20, $|Q_{20}^{+}|$ & 248 & 246 & 248 & & \\
Summariser calls & 382 & 382 & 261 & $-31.7$\,\% & $-31.7$\,\% \\
Mean radius $\bar k$ & whole document & 3 & 1.37 & & \\
Prompt volume (M characters) & 18.73 & 8.43 & 4.39 & $-76.6$\,\% & $-47.9$\,\% \\
Generation time, tracked scope & 54\,min~54\,s & 23\,min~48\,s\footnotemark[1] & 8\,min~45\,s & $-84.1$\,\% (6.3$\times$) & $-63.2$\,\% \\
Energy, tracked scope (Wh) & 106.68 & 52.52\footnotemark[1] & 13.22 & $-87.6$\,\% & $-74.8$\,\% \\
\quad of which GPU (Wh) & 46.78 & 21.73\footnotemark[1] & 8.80 & $-81.2$\,\% & $-59.5$\,\% \\
Emissions (\gco{}) & 21.78 & 10.72\footnotemark[1] & 2.83 & $-87.0$\,\% & $-73.6$\,\% \\
End-to-end preprocessing incl.\ PDF parsing & 71\,min~34\,s & 29\,min~09\,s & \todo{[not measured]} & & \\
Replicate run (same configuration) & --- & --- & 0.952 / 0.79791 / 3.92171; 543.7\,s; 13.51\,Wh; 2.86\,\gco{} & & \\
\botrule
\end{tabular}
\footnotetext[1]{Time and energy of the static $k=3$ row come from the tracked generation of 5 March 2026; its retrieval columns come from the ablation generation of the same configuration (a different summariser run; the conference paper's comparison used a third $k=3$ generation with \MRR{20} 0.77024 and \MAR{20} 4.01844 at the same \Rec{20}). CPU energy was measured by RAPL for the two March rows and estimated from load for the routed run; GPU energy is comparable across all rows. The tracked scope of the March $k=3$ row is to be confirmed (Sect.~\ref{sec:energy}).}
\footnotetext{\verify{Retrieval rows: \texttt{RnD/benchmarking\_retrievals.ipynb} (routed, baseline; $|Q_{20}^{+}|=250-2$ from the NaN counts printed there) and \texttt{RnD/ablation\_doc\_slice\_retrieval.ipynb} (static $k=3$; its NaN count 4 from \texttt{res\_ablation\_doc\_slice\_radius\_3.json} in \texttt{RnD/dynamic\_slice\_length/router\_experiments/results/}). Calls, $\bar k=(121\cdot0+261\cdot2)/382$: radius distribution printed by cell~2 of \texttt{RnD/ablation\_dynamic\_doc\_slice\_generation.ipynb}. Prompt volumes: \texttt{cost.py} re-implementation. Time, energy, GPU energy, emissions: rows 2026-03-05T23:59:43, 2026-03-05T22:53:27, 2026-09-25T20:06:47 of \texttt{RnD/emissions\_data/emissions.csv}; replicate: row 2026-09-25T16:06:28 and \texttt{RnD/dynamic\_slice\_length/router\_experiments/real\_runs/binary\_t2\_heldout.json}. End-to-end times: notebook comments cited in the text. Percentages: $1-\text{routed}/\text{reference}$; MRR and MAR deltas relative.}}
\end{table}

\begin{figure}[htbp]
  \centering
  \todo{[FIGURE INSERTION: grouped bar chart with three groups (full-document baseline, static $k=3$, routed pipeline) and three bars per group on two y-axes: total tracked energy in Wh (106.68, 52.52, 13.22), GPU energy in Wh (46.78, 21.73, 8.80) and emissions in \gco{} (21.78, 10.72, 2.83). Annotate the reductions relative to the baseline ($-50.8$\,\%, $-87.6$\,\% total; $-53.6$\,\%, $-81.2$\,\% GPU). Data: \texttt{RnD/emissions\_data/emissions.csv}, rows of 2026-03-05T23:59:43, 2026-03-05T22:53:27 and 2026-09-25T20:06:47. The conference figure \texttt{RnD/benchmarking\_results/preprocessing/eco\_friendliness.png} shows the first two groups.]}
  \caption{Tracked generation energy and greenhouse-gas emissions of the three pipelines on the 25-paper corpus. The static $k=3$ slice halves the baseline's energy and emissions; the routed pipeline removes 87.6\,\% of the energy and 87.0\,\% of the emissions. GPU energy, which was measured identically in every run, falls by 53.6\,\% and 81.2\,\% respectively.\verify{values and percentages as in Table~\ref{tab:primary}; $1-21.73/46.78=0.536$}}
  \label{fig:energy}
\end{figure}

\begin{figure}[htbp]
  \centering
  \todo{[FIGURE INSERTION: two-panel bar chart of preprocessing runtime. Panel (a) ``tracked generation scope'': 54.9, 23.8, 8.8\,min for baseline, static $k=3$, routed (rows of \texttt{emissions.csv} as in Fig.~\ref{fig:energy}). Panel (b) ``end-to-end incl.\ PDF parsing'': 71.6 and 29.2\,min for baseline and static $k=3$ (notebook comments; \texttt{RnD/benchmarking\_results/preprocessing/preprocessing\_time.png} plots these two as 4\,294 and 1\,749\,s) and a hatched placeholder bar for the routed pipeline marked ``\todo{[not measured]}''. Annotate speed-ups: 2.3$\times$ and 6.3$\times$ in (a), 2.5$\times$ in (b).]}
  \caption{Preprocessing runtime in the two scopes of Sect.~\ref{sec:energy}. In the tracked generation scope the static slice is 2.3 times and the routed pipeline 6.3 times faster than full-document contextualisation; in the end-to-end scope, which adds PDF parsing, the static slice is 2.5 times faster, and the routed pipeline has not yet been timed end to end.\verify{$3294.3/1428.2=2.31$, $3294.3/525.3=6.27$ (\texttt{duration} column of \texttt{RnD/emissions\_data/emissions.csv}); $4294/1749=2.46$ (notebook comments)}}
  \label{fig:runtime}
\end{figure}

\begin{figure}[htbp]
  \centering
  \todo{[FIGURE INSERTION: line chart, x-axis $K\in\{5,10,15,20\}$, y-axis \Rec{K}. Series: full-document baseline 0.73733 / 0.86500 / 0.93333 / 0.95400; routed pipeline 0.75167 / 0.86700 / 0.92600 / 0.95200; static $k=3$ 0.74100 / 0.85900 / 0.91067 / 0.94533; optionally no-context hybrid 0.71733 / 0.84700 / 0.91800 / 0.93933. Shade a $\pm0.7$-point band around the baseline at $K=20$ to show the single-run noise floor. Data: ``METRICS AT N='' blocks of \texttt{RnD/benchmarking\_retrievals.ipynb} and \texttt{RnD/ablation\_doc\_slice\_retrieval.ipynb}; the conference figures in \texttt{RnD/benchmarking\_results/retrieval/recall\_*.png} can be reused.]}
  \caption{Recall trajectories. The routed pipeline leads the full-document baseline at $K=5$ and $K=10$, trails it by 0.7 points at $K=15$ and by 0.2 points at $K=20$, inside the single-run band; the static $k=3$ slice trails at every cut-off beyond $K=5$.\verify{values listed in the placeholder; sources as in Table~\ref{tab:primary}}}
  \label{fig:recall}
\end{figure}

\subsection{Routing statistics, ablation of the routing policy and the noise floor}\label{sec:routing-results}

\subsubsection{What the router decided}
Of the 382 chunks, 42 were assigned $k=2$ by the table gate and 340 reached the SLM, which labelled 121 as self-describing ($k=0$, no summary) and 219 as needing their neighbours ($k=2$); the mean radius is $\bar k=1.37$, 31.7\,\% of the chunks were never sent to the summariser, and the 261 summaries that were written average 201 characters.\verify{radius and tier distribution printed by cell~2 of \texttt{RnD/ablation\_dynamic\_doc\_slice\_generation.ipynb} and the \texttt{doc\_slice\_radius}/\texttt{slice\_tier} fields of \texttt{RnD/preprocessed\_chunks/ablation\_doc\_slice\_radius\_dynamic.json} (0/slm 121, 2/slm 219, 2/table 42); $\bar k=522/382$; $121/382=0.317$; summary length: mean of \texttt{len(context)} over the 261 non-empty records} The replicate run made almost the same decisions (120 and 262).\verify{\texttt{k\_dist} of \texttt{RnD/dynamic\_slice\_length/router\_experiments/real\_runs/binary\_t2\_heldout.json}} The skipped chunks are what the prompt asks for: author and affiliation blocks, abstracts, conclusions, reference lists, acknowledgements and competing-interest statements, plus a few section stubs and short introductions.\verify{inspect the records with \texttt{doc\_slice\_radius == 0} in \texttt{RnD/preprocessed\_chunks/ablation\_doc\_slice\_radius\_dynamic.json}}

\subsubsection{How the policy was reached}
The routing policy of Sect.~\ref{sec:router} is the end point of an ablation that we summarise here; its full record, with more than forty simulated policies, is in the repository. The first router was a four-class genre classifier (whole-paper summary or boilerplate $\to$ no summary; self-contained exposition $\to k=1$; section-internal technical body $\to k=2$; fragment $\to k=3$) with tables gated to $k=3$. Its live run scored \Rec{20} of 0.9380, below fixed $k=2$ (0.94667) and even below the no-context control\verify{\texttt{RnD/router\_prompt\_engineering.md} Sect.~1 and the committed output of \texttt{RnD/ablation\_dynamic\_doc\_slice\_generation.ipynb} / \texttt{benchmarking\_retrievals.ipynb} at git commit a888f5c}; a per-question diagnosis showed that every gold chunk it lost had been routed to $k=2$ or was a table at $k=3$, and that the wide slices had produced degenerate table summaries (for one large table the $k=3$ summary was a verbatim copy of an example from the system prompt). Gating tables to $k=1$ recovered part of the loss (0.9420 in two identical live runs).\verify{\texttt{real\_runs/tables1.json} and \texttt{tables1\_run2.json} under \texttt{RnD/dynamic\_slice\_length/router\_experiments/}} An offline simulator that re-assembles a policy from stored summaries and replays the exact retrieval procedure then ranked the alternatives without summariser noise: $k=3$ is never needed once tables are gated below it; tables are best at $k=1$ or $2$ and worst at $3$; skipping self-describing chunks costs nothing for recall and gives the best MRR, MAR and \Rec{5}; narrowing technical body text to $k=1$ costs 0.9 points; and length-only or position-only heuristics (0.9373--0.9447) cannot replace the classifier.\verify{rows \texttt{sim\_len\_*} and \texttt{sim\_pos\_*} of \texttt{results/sweep\_results.jsonl} (Sect.~4 of \texttt{router\_prompt\_engineering.md})} A per-question oracle that picks the best $k\in\{1,2,3\}$ for every question reaches 0.964, so the headroom above fixed $k=2$ is about 1.7 points and only 17 of the 226 gold chunks are sensitive to $k$ at all.\verify{\texttt{scripts/oracle.py} over \texttt{results/res\_ablation\_doc\_slice\_radius\_1..3.json}; quoted in Sect.~3 of \texttt{router\_prompt\_engineering.md}} The simplest adequate policy, the two-class ``self-describing versus needs its neighbours'' prompt with tables at $k=2$, was the best offline (0.9507) and, once its exemplars were moved from benchmark papers to held-out papers, gave 0.9487 and then 0.952 in live runs.\verify{all numbers of this paragraph are in \texttt{RnD/router\_prompt\_engineering.md}: Sect.~1 (headline table), Sect.~2 (diagnosis of the 0.9380 run), Sect.~3 (oracle 0.964, 17 of 226 gold chunks $k$-sensitive), Sect.~4 (master table of every simulated policy, reproduced from \texttt{router\_experiments/results/sweep\_results.jsonl} by \texttt{scripts/make\_table.py}), Sect.~4b (round-3 modifications) and Sect.~8 (exemplar re-sourcing); live-run records in \texttt{router\_experiments/real\_runs/*.json}}

\subsubsection{Noise floor}
Regenerating the constant-$k=2$ configuration through the routed notebook reproduced its \Rec{20} of 0.94667 exactly, yet only 26 of the 380 summaries were identical to the earlier generation, 238 of the 250 questions received a different top-20 list (5 improved, 5 worsened) and \MRR{20} moved from 0.768 to 0.792.\verify{Sect.~3 of \texttt{RnD/router\_prompt\_engineering.md}; \texttt{real\_runs/fixed2.json} (0.94667, MRR 0.79227) versus the table \texttt{ablation\_doc\_slice\_radius\_2} (0.94667, MRR 0.76799); summary and top-20 comparison by \texttt{scripts/compare\_runs.py}} When, instead, the very same request sequence was issued twice from a freshly loaded model, all 380 summaries and every metric were identical.\verify{\texttt{tables1} versus \texttt{tables1\_run2} in \texttt{RnD/dynamic\_slice\_length/router\_experiments/real\_runs/} (identical metrics; Sect.~3 of the markdown report records 380/380 identical summaries)} The summariser's output therefore depends deterministically on the request history (the state of the server's cache and batching) rather than on a random seed, a single run resolves about 0.7 points of \Rec{20} and 2 points of \MRR{20}, and repeating an identical run does not sample this noise. Two consequences shape our reporting: policies are ranked offline on fixed summaries and confirmed live, and the two independent live generations of the final configuration (0.952 and 0.952) are the evidence that its result is not a favourable draw.\verify{the two generations are rows 2026-09-25T16:06:28 and 2026-09-25T20:06:47 of \texttt{emissions.csv}; metrics in \texttt{real\_runs/binary\_t2\_heldout.json} and in the saved output of \texttt{RnD/benchmarking\_retrievals.ipynb}} Moving the ten exemplars from benchmark papers to held-out papers changed the router's letter on 15\,\% of the chunks (mostly more skipping of front matter, acknowledgements and conclusions) but not a single question's outcome in the offline replay, which is why we regard the exemplar source as a matter of hygiene rather than of accuracy.\verify{Sect.~8 of \texttt{RnD/router\_prompt\_engineering.md}; letter agreement 287/338 between \texttt{results/routes\_binary.json} and \texttt{results/routes\_binary\_heldout.json} via \texttt{scripts/compare\_routes.py}}

\subsection{Retrieval quality and ranking dynamics}\label{sec:ranking}

The ranking metrics explain why skipping a third of the summaries does not hurt. A summary is prepended to the chunk before embedding, so for a chunk that already states what it is about (an abstract, a conclusion) the summary only dilutes the representation; embedding such chunks raw is what lifts the routed pipeline's \Rec{5} above the baseline (0.75167 against 0.73733) and its \MRR{20} to 0.79798, the highest of all configurations in this study.\verify{``METRICS AT N=5'' and ``N=20'' blocks of \texttt{RnD/benchmarking\_retrievals.ipynb}; compare the \MRR{20} column of Table~\ref{tab:sweep}} With \MAR{20} of 3.91 the gold chunks of a question sit on average at rank four of twenty, i.e.\ in the first fifth of the context that a downstream generator would receive, which is the primacy region of the U-shaped accessibility curve of Fig.~\ref{fig:ucurve}; the baseline (3.94) and the static slice (4.02) are marginally deeper. The number of questions for which nothing relevant is retrieved is small and similar across systems (17, 4, 3 and 2 of 250 at $K=5,10,15,20$ for the routed pipeline against 17, 6, 2 and 2 for the baseline), so the ranking averages are computed over essentially the same question sets.\verify{\Rec{5}, MRR, MAR and NaN counts from the printouts of \texttt{RnD/benchmarking\_retrievals.ipynb}; baseline NaN counts 17/6/2/2 from \texttt{res\_anthropic\_control\_table.json} in \texttt{RnD/dynamic\_slice\_length/router\_experiments/results/}}

\begin{figure}[htbp]
  \centering
  \todo{[FIGURE INSERTION: line chart of retrieval accuracy versus the position of the relevant passage in the context, redrawn schematically after Liu et al.\ \cite{liu2023lostmiddlelanguagemodels} (U-shape: high at the first positions, minimum in the middle, partial recovery at the end; the conference paper used \texttt{images/u\_curve\_no\_border.png}, not in the repository). Overlay, on a 20-position axis, markers at the mean average rank of the three systems: 3.94 (full document), 4.02 (static $k=3$), 3.91 (routed), and shade positions 1--4 as the ``top 20\,\%'' region.]}
  \caption{Positional accessibility of the retrieved evidence. Language models use information at the beginning and end of a prompt better than in its middle; with a mean average rank of about four in a 20-chunk context, all three pipelines place the supporting chunks in the primacy region, the routed pipeline marginally earliest.\verify{MAR values from Table~\ref{tab:primary}}}
  \label{fig:ucurve}
\end{figure}

\subsection{Cross-domain robustness on European Commission policy documents}\label{sec:policy}

Table~\ref{tab:policy} repeats the static comparison on the policy corpus. The fixed $k=3$ slice reaches \Rec{20} of 0.9625 against 0.96875 for the full-document baseline ($-0.6$ points), \MRR{20} 0.813 against 0.822 and \MAR{20} 3.62 against 3.52, while its tracked generation took 252\,s instead of 858\,s ($-70.6$\,\%) and consumed 5.08 instead of 18.00\,Wh ($-71.8$\,\%; GPU energy $-72.1$\,\%).\verify{metrics: the ``METRICS AT N='' printout for the 80-question set saved in the committed version of \texttt{RnD/benchmarking\_retrievals.ipynb} of 21 September 2026 (git commit f227c29) and \texttt{RnD/benchmarking\_results/tobacco/dashboard\_view\_recall\_tobacco.png}; time and energy: rows 2026-07-19T15:01:49 (\texttt{tobacco\_document\_slice}) and 2026-07-19T16:44:18 (\texttt{tobacco\_traditional}) of \texttt{RnD/emissions\_data/emissions.csv}} The recall gap is the same size as on the scientific corpus, and the efficiency gain is larger because the policy documents are longer (up to 60 chunks, 103\,000 characters), which is exactly the regime in which Proposition~\ref{prop:complexity} predicts the widest gap. Regulatory text, with its numbered articles, member-state tables and annexes, is thus served by the bounded slice at least as well as scientific prose. The routing agent has not yet been run on this corpus; its column is left open, and the prompt of Listing~\ref{lst:router}, which names abstracts and conclusions, will need policy-specific exemplars (executive summaries, legal bases, annex tables) before the comparison is meaningful.

\begin{table}[htbp]
\caption{Out-of-distribution evaluation on the European Commission policy corpus (80 multi-hop questions, 4 documents, 141 chunks).}\label{tab:policy}
\begin{tabular}{@{}lcccc@{}}
\toprule
Metric & Full document & Static $k=3$ & Routed pipeline & $\Delta$ $k=3$ vs full \\
\midrule
\Rec{5} & 0.78958 & 0.75833 & \todo{[run]} & $-3.1$\,pt \\
\Rec{10} & 0.90625 & 0.90000 & \todo{[run]} & $-0.6$\,pt \\
\Rec{15} & 0.95000 & 0.94375 & \todo{[run]} & $-0.6$\,pt \\
\Rec{20} & 0.96875 & 0.96250 & \todo{[run]} & $-0.6$\,pt \\
\MRR{20} & 0.82218 & 0.81296 & \todo{[run]} & $-1.1$\,\% \\
\MAR{20} & 3.51667 & 3.61875 & \todo{[run]} & $+2.9$\,\% \\
Generation time, tracked scope & 14\,min~18\,s & 4\,min~12\,s & \todo{[run]} & $-70.6$\,\% \\
Energy, tracked scope (Wh) & 18.00 & 5.08 & \todo{[run]} & $-71.8$\,\% \\
\quad of which GPU (Wh) & 14.53 & 4.05 & \todo{[run]} & $-72.1$\,\% \\
Emissions (\gco{}) & 3.68 & 1.04 & \todo{[run]} & $-71.8$\,\% \\
\botrule
\end{tabular}
\footnotetext{\todo{[confirm two properties of the stored policy tables before publication: (i) the document-slice table (\texttt{tobacco\_sliding\_table}) was built with the raw chunk text as the embedding source, so its summaries reached only the BM25 index, whereas the baseline table embeds summary plus chunk; (ii) whether the tracked $k=3$ run wrapped the PDF conversion. Either would make the slice column conservative.]} \verify{metrics as cited in the text; energy components from the two July rows of \texttt{RnD/emissions\_data/emissions.csv} (\texttt{energy\_consumed}, \texttt{gpu\_energy}, \texttt{emissions}, \texttt{duration}); deltas $=1-k3/\text{full}$ for time and energy, relative for MRR and MAR, absolute for recall}}
\end{table}

\subsection{Hardware profiling}\label{sec:profiling}

The memory model of Sect.~\ref{sec:complexity} was not instrumented per run, and Fig.~\ref{fig:memory-surface} is therefore a placeholder for the measurement we consider necessary. What the repository records is consistent with the model: a peak of 7\,695\,MiB of device memory over the whole pipeline, a resident footprint of about 4.6\,GB for the summariser runner during slice generation, and, from the energy tracker, a mean GPU utilisation of 46.5\,\% for the full-document baseline against 90.1\,\% for the routed pipeline and a mean host-memory occupancy of 15.9\,GB against 17.3\,GB.\verify{\texttt{README.md} line~59 (7\,695\,MiB); cell~4 comment of \texttt{RnD/doc\_slice\_chunking.ipynb} (4\,644\,MB); \texttt{gpu\_utilization\_percent} and \texttt{ram\_used\_gb} of the rows timestamped 2026-03-05T23:59:43 and 2026-09-25T20:06:47 of \texttt{RnD/emissions\_data/emissions.csv}} The low average utilisation of the baseline, at a higher reserved context window and a per-call prompt of up to 18\,000 tokens, is what a bandwidth-bound run looks like from the outside; the conference study additionally observed host-memory growth approaching 10\,GB and a resident device footprint of about 5\,GB for the slice pipeline, but those readings were taken by hand and are marked for re-measurement rather than reported as data.

\begin{figure}[htbp]
  \centering
  \todo{[FIGURE INSERTION: 3-D surface or scatter plot. Axes: reserved context window (tokens; at least the two operating points 16\,384 and 30\,000), document size (chunks, 8--33 for the scientific corpus, up to 60 for the policy corpus) and resident device memory (GB) with host-memory overflow shown as a second surface or colour. Mark the 8\,GB device limit as a plane. All points are \todo{[to be measured with nvidia-smi and psutil at 1-s intervals during a baseline run and a routed run]}; the only firm values are the 7\,695\,MiB peak and the mean host memory in use of 15.9\,GB (baseline) and 17.3\,GB (routed) from the energy tracker.]}
  \caption{Memory footprint versus reserved context window and document size (placeholder). The bounded slice keeps the reserved window at 16\,384 tokens for every document and stays under the 8\,GB device limit; the full-document baseline reserves 30\,000 tokens and crosses it, paging to host memory. Measured points are to be added.}
  \label{fig:memory-surface}
\end{figure}

\section{Discussion and future trajectories}\label{sec:discussion}

\subsection{Operational calculus}
For an authority indexing thousands of pages on premises, the relevant frontier is accuracy against time and energy. On that frontier the full-document baseline buys its last 0.2 points of \Rec{20} with 8.1 times the generation energy and 6.3 times the wall-clock of the routed pipeline, and the static $k=3$ slice sits in between, at half the baseline's energy for a 0.9-point loss.\verify{$106.68/13.22=8.07$ and $3294.3/525.3=6.27$ from Table~\ref{tab:primary}; $0.954-0.94533=0.0087$} Because every recall gap between the three systems lies inside the $\pm0.7$-point band of a single run, whereas the energy gaps exceed any run-to-run variation by an order of magnitude, the Pareto choice is unambiguous: bounded, routed contextualisation dominates. The absolute quantities are small, but they scale with corpus size, recur at every re-index, and the baseline's paging penalty grows with document length.

\subsection{Ablation synthesis}
Replacing the static heuristic with the routing agent gained 0.7 points of \Rec{20} and removed 63\,\% of the generation time and 75\,\% of the energy of the $k=3$ slice. The development path (0.938, 0.942, 0.9487, 0.952) shows that the gain came from the class-to-radius mapping, tables never at $k=3$ and self-describing chunks at $k=0$, rather than from prompt wording, and that $k=3$ is unnecessary once tables are gated. The router costs one token per chunk, about a minute per corpus and no second model. Its limits are equally clear: the oracle headroom above fixed $k=2$ is 1.7 points and the summariser's run-to-run variation is of the same order, so further gains must come from the summariser or from multi-run evaluation, not from finer routing classes.\verify{path and headroom: Sects.~1, 3 and 8 of \texttt{RnD/router\_prompt\_engineering.md}; $0.964-0.94667=0.017$}

\subsection{Architectural constraints and roadmap}
The slice is intra-document: a chunk whose meaning depends on another document, an annex cited in a report or a standard referenced by an audit, receives no help. Thousand-page audit or financial filings have not been tested; Proposition~\ref{prop:complexity} predicts the largest gains there, but also the largest exposure to cross-section dependencies. The natural extension is edge-to-cloud collaborative routing: the agent already exposes a confidence margin, so near-tie chunks, or questions a local generator cannot answer from the retrieved context, can be escalated selectively to a larger model while low-entropy work stays local. Next steps are a router run on the policy corpus with policy-specific exemplars, per-run memory instrumentation, and enough repeated generations to replace single-run bands by interval estimates.

\section{Conclusion}\label{sec:conclusion}

Index-time contextual retrieval can be made to fit the memory, time and energy budget of a single consumer GPU without giving up its accuracy. Bounding each chunk's context to a window of neighbours makes prompt length and reserved memory independent of document length, and a one-token routing decision of the same small model removes a third of the summaries and most of the remaining cost. On 250 multi-hop questions over 25 scientific papers the routed pipeline matched full-document contextualisation within the noise of a single run (\Rec{20} 0.952 against 0.954) with better ranking metrics, at 12\,\% of its generation energy and 16\,\% of its generation time, and the fixed-radius variant transferred to European Commission policy documents with the same accuracy gap and a 72\,\% energy saving.\verify{$13.22/106.68=0.124$ and $525.3/3294.3=0.159$ (Table~\ref{tab:primary}); $1-5.08/18.00=0.718$ (Table~\ref{tab:policy})} The pipeline is precise (its metrics are reproducible per question), explainable (every routing decision is logged with its evidence) and provable (every run carries an energy and emission ledger), which is the PEP-AI footing we consider necessary for AI-assisted compliance auditing at the edge of the energy system.

\backmatter

\bmhead{Acknowledgements}
An earlier version of this work was presented at the Student Research Conference of the University of Debrecen and published as a conference paper \cite{sandor2025documentslicecontextualretrieval}. It was supported by the UD Talent Program and by the Doctoral School of Informatics and the Faculty of Informatics of the University of Debrecen.

The authors acknowledge the use of generative artificial intelligence tools, specifically Gemini (Google), Claude (Anthropic), ChatGPT (OpenAI) and Grammarly, in preparing this manuscript and the associated research. These systems were used to improve the stylistic quality, grammatical accuracy and coherence of the text; Gemini assisted the initial literature survey and generated the silver-standard question sets under the protocol of Sect.~\ref{sec:datasets}; Gemini and ChatGPT were used as coding assistants; and Claude was used to draft this journal version from the conference manuscript and the repository's experimental records, including the bibliography entries marked as machine-added in the source file, and to run the dynamic-radius routing experiments under the authors' direction. The core concepts, research methodology and scientific contributions are the authors' original work; all AI-generated outputs, including text, source code and the experimental protocol, were produced under human supervision, were validated manually and were critically reviewed and edited by the authors, who remain fully responsible for the accuracy, integrity and originality of the content.

\section{Declarations}

\begin{itemize}
\item \textbf{Funding:} UD Talent Program; Doctoral School of Informatics and Faculty of Informatics, University of Debrecen \todo{[confirm grant identifiers, if any]}.
\item \textbf{Conflict of interest:} The authors declare no competing interests.
\item \textbf{Ethics approval and consent to participate:} Not applicable.
\item \textbf{Consent for publication:} Not applicable.
\item \textbf{Data availability:} The frozen chunk files, question sets, per-run energy records (\texttt{emissions.csv}), routing records and per-question retrieval results analysed in this article are available in the project repository \cite{sandor2025documentslicecontextualretrieval}.
\item \textbf{Materials availability:} Not applicable.
\item \textbf{Code availability:} The complete pipeline (parsing, routing agent, slice generation, indexing, evaluation and the offline policy simulator) is available in the same repository \cite{sandor2025documentslicecontextualretrieval}.
\item \textbf{Author contribution:} \todo{[insert CRediT statement]}.
\end{itemize}

\bibliography{refs}

\end{document}
